
%% ================================================================%%
%% thesis.tex - template for an honors thesis in mathematics       %%
%% ================================================================%%
%% Save this page to a file named ``thesis.tex''                     %%
%% ================================================================%%

\documentclass[11pt,twoside]{report}

% If you want to print single-sided, then to make the margins work
% correctly, change the above line to
%          \documentclass[11pt]{report}
% (without the %, of course)



% Import standard math macro packages

\usepackage{hyperref}
\hypersetup{
    colorlinks=true,
    linkcolor=blue,
    filecolor=blue,      
    urlcolor=blue,
}

\urlstyle{same}
\usepackage{amsmath}
\usepackage{amsthm}
\usepackage{amssymb}
\usepackage{amsfonts}
\usepackage{graphicx}
\usepackage{bbm}
\usepackage{esvect}
\usepackage{listings}
\usepackage{wrapfig}
\usepackage{tikz}

% Declare theorem types.  Add, subtract, and modify
% as you wish.

\newtheorem{thm}{Theorem}
\newtheorem{lem}[thm]{Lemma}
\newtheorem{prop}[thm]{Proposition}
\newtheorem{cor}[thm]{Corollary}
\newtheorem{dfn}[thm]{Definition}
\newtheorem{definition}[thm]{Definition}
\newtheorem*{conj}{Conjecture}
\newtheorem{obs}[thm]{Observation}
\newtheorem{ex}{Exercise}[section]
\newtheorem{example}[thm]{Example}
\newtheorem{fact}[thm]{Fact}
\newtheorem{claim}[thm]{Claim}
\newtheorem{cla}[thm]{Claim}
\newcommand{\w}{\omega}
\DeclareMathOperator*{\argmax}{argmax}
\numberwithin{thm}{section}
\numberwithin{equation}{section}


% Set margins and spacing.  Probably you shouldn't mess with these,
% unless your printer isn't aligned quite the same as mine and the
% margins are coming out wrong.  (At least 1.5 inches on binding side,
% at least 1 inch on the other three sides.)

\setlength{\evensidemargin}{0in}
\setlength{\oddsidemargin}{0.5in}
\setlength{\textwidth}{6in}
\setlength{\topmargin}{0.2in}
\setlength{\textheight}{8.6in}
\setlength{\footnotesep}{14pt}


%%%%%%%%%%%%%%%%%%%%%%%%%%%%%%%%%%%%%%%%%%%%%%%%%%%%%%%%
% IF YOUR PRINTER ISN'T CENTERING THE PAGES CORRECTLY: %
%%%%%%%%%%%%%%%%%%%%%%%%%%%%%%%%%%%%%%%%%%%%%%%%%%%%%%%%
% the following lines adjust the vertical and horizontal
% positioning of the entire page.
% (As noted above, you want different margins on the right
% and left, so when I say ``centering'' above, it
% technically isn't horizonally centered.)
% \voffset is how much to slide the text down the page.
% \hoffset is how much to slide the text to the right on the page.
% Either quantity can be negative.

\setlength{\voffset}{-0.7in}
\setlength{\hoffset}{0in}


% User-defined macros.  Add, subtract, modify as you wish.

\newcommand{\C}{{\mathbb{C}}}
\newcommand{\F}{{\mathbb{F}}}
\newcommand{\N}{{\mathbb{N}}}
\newcommand{\Q}{{\mathbb{Q}}}
\newcommand{\PP}{{\mathbb{P}}}
\newcommand{\R}{{\mathbb{R}}}
\newcommand{\Z}{{\mathbb{Z}}}

\newcommand{\cK}{{\mathcal{K}}}

\newcommand{\fm}{{\mathfrak{m}}}
\newcommand{\fo}{{\mathfrak{o}}}

\newcommand{\Cp}{\C_p}
\newcommand{\Cv}{\C_v}
\newcommand{\Qp}{\Q_p}
\newcommand{\Qpbar}{\bar{\Q}_p}
\newcommand{\Dbar}{\overline{D}}
\newcommand{\Rd}{\mathbb{R}^{d}}
\newcommand{\dx}{\mathrm{d}x}
\newcommand{\la}{\left<}
\newcommand{\ra}{\right>}


% User-defined math operators.  Add, subtract, modify as you wish.

\DeclareMathOperator{\rad}{rad}
\DeclareMathOperator{\diam}{diam}
\DeclareMathOperator{\divop}{div}
\DeclareMathOperator*{\argmax}{arg\,max}


%%%%%%%%%%%%%%%%%%%%%
% Let's get started %
%%%%%%%%%%%%%%%%%%%%%

\begin{document}


% Set up double-spacing (must be done *after* \begin{document})

\setlength{\baselineskip}{21pt}

% Roman numeral page numbers for abstract, etc.

\pagenumbering{roman}


%%%%%%%%%%%%%%
% Title page %
%%%%%%%%%%%%%%

\begin{titlepage}
\mbox{}
\vspace{1in}
\begin{center}
\LARGE CleanFM: A Windowless Approach to Spectral Leakage and the Frequency Modulation Transform \\
 [2in]
\normalsize  Matthew McQuistion \\[\medskipamount]
April 25th, 2025 \\[1in]
Submitted to the \\
Department of Applied Mathematics \\
of Brown University  \\

\end{center}
\vfill
\begin{center}
Faculty Advisor: Professor Amalia Culiuc
\end{center}
\vfill
\begin{center}
Copyright \copyright\ 2025 Matthew McQuistion
\end{center}
\end{titlepage}


%%%%%%%%%%%%



%%%%%%%%%%%%%%%%%%%%
% Acknowledgements %
%%%%%%%%%%%%%%%%%%%%

\chapter*{Acknowledgments}
\addcontentsline{toc}{chapter}{\numberline{}Acknowledgments}

I would like to express my sincere gratitude to my thesis advisor, Professor Amalia Culiuc, for her invaluable guidance and support throughout the last year. I am also thankful for my friends and professors in the Applied Mathematics and Computer Science departments at Brown University along with their insights and discussions. This thesis would not have been possible without them.

The source code for all algorithms presented in this thesis is publicly available at \\\url{https://github.com/matthewmcq/cleanfm} and licensed under Apache 2.0.

%%%%%%%%%%%%%%%%%%%%%
% Table of contents %
%%%%%%%%%%%%%%%%%%%%%



\tableofcontents


%%%%%%%%%%%%%%%%%%%%%%%%%%%%%%%%%%%%%%
% If you also have a lot of figures
% or tables, un-comment one or both
% of the following lines to generate
% a list of them.
%%%%%%%%%%%%%%%%%%%%%%%%%%%%%%%%%%%%%%

\listoffigures
%\listoftables


%%%%%%%%%%%%%%%%%%%%
% List of notation %
%%%%%%%%%%%%%%%%%%%%

\chapter*{List of Notation}
An important note to any readers with more of a mathematical background than one of signal processing is that certain common words have entirely different meanings across the two disciplines.

Specifically, the ``spectrum'' of a signal throughout the context of this thesis refers to the set of component frequencies obtained by computing its Fourier Transform. Similarly, ``spectral'' does not involve the Spectral Theorem nor does it refer to eigenvalues in any direct way. Furthermore, the term ``harmonic'' 
 refers to any integer-multiples of a lower (fundamental) frequency. 

\begin{itemize}

    \item \textbf{Signal:} In the context of this thesis, a signal refers to any time-varying quantity that can be mathematically represented as a function of time, typically denoted as $f(t)$ or $x(t)$ for continuous signals, or as a sequence $x[n]$ for discrete signals. For audio applications, signals represent physical quantities such as air pressure variations (sound waves) or their electrical/digital analogues. More broadly, a signal can be any measurable phenomenon that varies over time, including but not limited to audio waveforms, electromagnetic waves, stock prices, or biological measurements.

    \item \textbf{WAV File (.wav): }The Waveform Audio File Format is a lossless file format used for high quality audio data. Audio is encoded into the WAV format through techniques such as pulse code modulation and directly stores a representation of the relative sound pressure level at each point in time, with each data point corresponding to a single sample. 

    \item \textbf{Frequency: }The frequency of a sinusoid corresponds to $\frac{1}{\lambda}$ where $\lambda$ is its wavelength. In a musical context, the frequency of a note corresponds to its pitch. 

    \item \textbf{Filter: }A filter is a function that modifies the frequency spectrum of a signal. The most common examples are high and low-pass filters, which remove all frequencies below or above, respectively, the choice of some other frequency, which is called the filter's ``cutoff.'' 

    \item \textbf{Impulse Response: }The impulse response of a system refers to its behavior after a brief, intense signal. We are typically concerned with the impulse responses of filters, but the term can apply more broadly to other mathematical objects, such as kernels. 

    \item \textbf{ADSR Envelope: }The ADSR envelope refers to the shape of a note's waveform, which is given by its volume over time. The four parts of an ADSR envelope are:
    \begin{enumerate}
        \item \textbf{A}ttack: The time it takes from the start of the impulse to the point at which its amplitude is highest. 
        \item \textbf{D}ecay: The time it takes to get from the initial maximum amplitude peak to a sustained volume.
        \item \textbf{S}ustain: The volume at which the note will remain if constantly held. 
        \item \textbf{R}elease: The time it takes for all sound to stop after the note ceases to be activated. 
    \end{enumerate}

    \item \textbf{Sample Rate: }The sample rate of a signal is the frequency at which  samples are taken or measured. Common sample rates for high-quality audio data are 44.1kHz and 48kHz. 

    \item \textbf{Nyquist Frequency: }The Nyquist frequency/rate/limit of a signal is the highest frequency that can appear undistorted and without aliasing. It is equal to exactly one half of the sample rate.\footnote{See Appendix \ref{sec:NQST}} The sample rates of 48kHz and 44.1kHz are used because they have a Nyquist frequency above 20kHz, which is the threshold for human hearing. Also, the DFT of a signal has Hermitian (conjugate) symmetry about the Nyquist frequency. 

    \item \textbf{Aliasing: }Aliasing occurs when signals are sampled below their Nyquist rate, which can cause high-frequency components to ``fold back'' into incorrect lower frequencies, distorting the signal and introducing undesirable artifacts. 

    \item \textbf{Resampling: }Resampling (\ref{sec:resample}) is the process of changing the sample rate of a signal. If the new sample rate is lower, the process is also called ``downsampling'' and can introduce aliasing if the signal has frequencies below the old Nyquist limit but above the new one. To avoid this, a low-pass filter is commonly used before downsampling to remove these frequencies. If the new sample rate is higher, the process is called ``upsampling'' and is less likely to introduce these types of artifacts, but will not introduce any new frequency content above the old sample rate's Nyquist frequency. 

    \item \textbf{Short-Time Fourier Transform (Spectrogram): }The Short-Time Fourier Transform (\ref{subsec:cstft_intro}) computes individual DFT slices of a longer audio segment to give a time-frequency representation of the signal that would not be possible with a single DFT. The STFT ``Spectrogram'' refers to a plot of the magnitudes of the STFT over time, where the x-axis corresponds to time, y-axis to frequency, and the pixel intensity to magnitude. Each vertical slice corresponds to one individual DFT. 

    \item \textbf{Analog-to-Digital Converter (ADC): }A system that converts an analog (continuous, real-world) signal into a digital (discrete, sampled) signal. A digital-to-analog converter (DAC) does the opposite.

    \item \textbf{(FM) Operator: }An FM operator refers to one of the simple waveforms that comprises an FM synthesizer. These include the carrier and individual modulators. Note that typically an FM operator includes the oscillator along with an envelope generator and an amplifier, for our purposes we will not consider the envelope generator, as they have complex mathematical interactions with the frequency spectrum that go outside the scope of this paper. Also, as our exploration of FM is strictly mathematical, we do not need to worry about the amplifier. 

\end{itemize}


%%%%%%%%%%%%%%%%%%%%%%%%%%
% Finish up front matter %
%%%%%%%%%%%%%%%%%%%%%%%%%%

% Some magic code to leave an extra blank page
% if we are double-siding AND the front matter
% ends on an odd-numbered page.
% (Otherwise the even- and odd-side margins
% will be backwards for the rest of the thesis.)

\makeatletter 
\if@twoside \ifodd\value{page}
\clearpage\mbox{}\thispagestyle{empty} \fi \fi
\makeatother 


% Switch back to arabic page numbers AFTER the end of the current page.

\clearpage
\pagenumbering{arabic}

\chapter*{Introduction}
At the heart of modern digital signal processing lies the fundamental challenge of accurately representing continuous real-world phenomena in a finite number of ones and zeros. Since the advent of the Fast Fourier Transform (FFT) in the 1960s, we have relied on frequency-domain analysis to understand, manipulate, and transform signals ranging from audio recordings to biomedical measurements. However, the discretization inherent in these methods introduces artifacts and limitations that have shaped—and constrained—the field for decades.

One of the most persistent challenges is that of spectral leakage, where energy from a single frequency component ``spills'' into adjacent frequency bins, creating a smearing effect that obscures the true spectral content. Traditionally, this problem has been addressed through window functions that attenuate the edges of analysis frames at the cost of frequency precision, amplitude accuracy, and signal resolution. This tradeoff between leakage reduction and resolution has become an accepted compromise in signal processing pipelines.

This thesis presents a fundamentally different approach. Rather than accepting spectral leakage as an inevitable consequence of discretization, we introduce mathematical techniques that directly confront and resolve it by working with the underlying signal characteristics.

First, we develop the method of Dirichlet Kernel Deconvolution (DKD), which extracts continuous frequency components from discrete spectra by modeling and removing the intrinsic spectral leakage pattern. Unlike traditional approaches that apply windows before transformation, DKD operates on the already-transformed spectrum, preserving both the full resolution and complete information content of the signal. We then extend this method with an efficient parallel implementation that maintains accuracy while dramatically reducing computational requirements.

Building on these foundations, we introduce the Frequency Modulation Transform (FMT), a novel mathematical operator that reveals harmonic relationships within complex signals by analyzing the ``frequencies of frequencies.'' The FMT unlocks new capabilities for timbral analysis and source separation that previously required sophisticated machine learning techniques.

Finally, we demonstrate practical applications of these methods, including an idealized approach to signal resampling, a ``cleaned'' version of the Short-Time Fourier Transform with superior frequency resolution, and a technique for source separation that requires no training data yet achieves compelling preliminary results.

Together, these contributions offer a new perspective on spectral analysis that challenges longstanding assumptions about the inevitable tradeoffs in digital signal processing. By addressing these fundamental limitations at their mathematical core, we open possibilities for more accurate, efficient, and insightful signal processing across diverse applications.

The remainder of this thesis is organized as follows: Chapter \ref{ch:1} provides background on the Discrete Fourier Transform, timbre, and frequency modulation. Chapter \ref{ch:3} provides motivation for the origins of spectral leakage and window functions along with a brief overview of current approaches. Chapter \ref{ch:4} introduces the DKD algorithm in both single-threaded and parallel forms. Chapter \ref{ch:5} discusses applications of DKD such as resampling and ``cleaning'' the STFT. In Chapter \ref{ch:2}, we pivot and explore both FM and AM Synthesis and their mathematical properties. Finally, Chapters \ref{ch:6} and \ref{ch:7} present the Frequency Modulation Transform and its applications to audio processing. 
\chapter{Background}
\label{ch:1}
\section{The Discrete Fourier Transform}
Fourier analysis is a ubiquitous and powerful tool in mathematics, and at its root lies the Fourier series: a linear combination of sine and cosine functions of different frequencies that converges to another (periodic) function defined in rectangular coordinates~\cite{Lighthill_1958}.

Frequency, in this context, refers to the number of complete cycles a sinusoidal component makes per unit time. For example, a sine wave with frequency 440Hz completes 440 cycles per second. The general form of a Fourier series represents a periodic function $f(t)$ with period $T$ as:
\begin{equation}
  f(t) = \frac{a_0}{2} + \sum_{n=1}^{\infty} \left[ a_n \cos\left(\frac{2\pi nt}{T}\right) + b_n \sin\left(\frac{2\pi nt}{T}\right) \right]  
\end{equation}

where the coefficients $a_n$ and $b_n$ determine the amplitude of each frequency component.

The operation converting the original function into its Fourier series is called the (continuous) Fourier transform, $\mathcal{F}$, defined for all angular frequencies $\omega \in \mathbb{R}$ as: 
\begin{equation}
  \mathcal{F}(\omega) = \int_{-\infty}^\infty f(t) e^{-i\omega t}dt  
\end{equation}
where the angular frequency $\omega = 2\pi\xi$ for some $\xi \in \mathbb{R}$. This transform is well-defined for any absolutely integrable function ($f \in L^1(\mathbb{R})$), meaning functions where $\int_{-\infty}^{\infty} |f(t)| dt < \infty$. The domain can be extended to include a larger class of functions, such as square-integrable functions ($f \in L^2(\mathbb{R})$), through appropriate mathematical techniques. For practical digital signal processing applications, we work with bounded, finite-duration signals that satisfy these conditions.

 In other words, the Fourier transform evaluated for a particular frequency $\omega$ can be interpreted as the convolution of the original function with the sine and cosine functions of that particular frequency. This comes from Euler's identity, which states that:
 \begin{equation}
  e^{i\theta} = \cos\theta + i\sin\theta   
 \end{equation}
 One of the most important properties of the Fourier transform is that it is invertible, meaning that one can apply an inverse Fourier transform to the Fourier series and obtain the original function.

The inverse Fourier transform is given by:

\begin{equation}
  f(t) = \frac{1}{2\pi}\int_{-\infty}^{\infty} \mathcal{F}(\omega) e^{i\omega t} d\omega  
\end{equation}

This pair of transformations allows us to move seamlessly between the time domain (where signals are represented as functions of time) and the frequency domain (where signals are represented by their frequency components).

Another fundamental property of the Fourier transform is that convolution in the time domain corresponds to multiplication in the frequency domain~\cite{BLACKLEDGE200530}. For two functions $f(t)$ and $g(t)$ with Fourier transforms $F(\omega)$ and $G(\omega)$ respectively, we have:
\begin{equation}
\label{eq:conv}
  \mathcal{F}\{f \ast g\} = F(\omega) \cdot G(\omega)  
\end{equation}

where $\ast$ represents the convolution operation. This property is particularly valuable in signal processing, as it is sometimes more efficient to perform convolution operations (such as filtering) by converting to the frequency domain, performing simple point-wise multiplication, and then converting back.

Although this method of computing the Fourier transform is incredibly useful throughout many branches of mathematics, modern computers cannot accurately represent continuous data. In order to analyze real-world signals using a computer, we must first discretize the data into a finite number of samples that approximate the continuous signal. Furthermore, this process of sampling must follow the Nyquist-Shannon sampling theorem (see Appendix \ref{sec:NQST}), which states that to perfectly reconstruct a continuous, bandlimited signal, we must sample at a rate greater than twice the highest frequency component in the signal~\cite{Shannon1949}.

However, even with proper sampling, we still face an issue: to use the Fourier transform above, we would need our data to be in a continuous interval, not a finite collection of points. The solution to this problem is to use the Discrete Fourier Transform, which converts an $N$ point sequence of complex numbers $\{\mathbf{x}_n\} := x_0, x_1, \dots, x_{N-1}$ into a new sequence of complex numbers $\{\mathbf{X}_k\} := X_0, X_1, \dots, X_{N-1}$ defined by:
\begin{equation}
  X_k = \sum_{n=0}^{N-1}x_n \cdot e^{-i2\pi\frac{k}{N}n}  
\end{equation}

\subsection{Convergence of the Discrete Fourier Transform}

When dealing with digital signal processing (DSP), we work with the Discrete Fourier Transform (DFT) rather than the continuous version. An important advantage of the DFT is that convergence is guaranteed for any finite discrete signal~\cite{Lai-Man_2022}, which encompasses virtually all practical DSP applications.

The mathematical reason for this guaranteed convergence is that the DFT operates on finite-length sequences and involves only finite summations rather than infinite integrals. Specifically, for a sequence of $N$ complex numbers, the DFT formula shown above always produces a well-defined result.

Unlike the continuous case where issues of convergence can arise with certain pathological functions, the discrete nature of digital signals ensures that the DFT always converges. This mathematical certainty is one reason why the DFT has become such a foundational tool in digital signal processing, audio analysis, and numerous other computational applications.

In addition to being invertible and linear, the DFT behaves nearly identically to its continuous analogue with a few notable exceptions~\cite{Benson_2006}. These exceptions include:

\begin{itemize}
    \item \textbf{Spectral Leakage}: When the signal's period doesn't match the DFT interval, energy from one frequency ``leaks'' into surrounding frequency bins.
    
    \item \textbf{Aliasing}: When signals contain frequencies above the Nyquist limit, these components appear as lower-frequency artifacts in the DFT output.
    
    \item \textbf{Resolution Limitations}: The DFT can only resolve frequencies at discrete intervals determined by the sample rate and window length.
\end{itemize}

Despite these limitations, the DFT can be computed with $\mathcal{O}(n\log n)$ time complexity thanks to the aptly named ``Fast Fourier Transform'' (FFT), making it computationally efficient for real-world applications.





\section{Fundamentals, Partials, and Timbre}
``Concert Pitch'' is the note to which instruments are tuned in an orchestral setting. The ISO defines this pitch to be A440, or A above middle C (A4), which has a frequency of exactly 440.0Hz. Despite the fact that all instruments are tuned to match this, the actual sound that a note at this pitch makes varies substantially between instruments. While the fundamental frequency, the lowest frequency of a periodic waveform, is 440Hz, the sound that a piano produces when A above middle C is played is far richer than that of a simple isolated sine wave corresponding to that frequency. The reason for this is because of timbre, an intrinsic property of the instrument, which describes the sonic characteristics of a note beyond just its pitch. Essentially, there is a fundamental ``piano-ness'' to any note played on a piano which uniquely identifies it as a sound produced by a piano rather than one produced by a different source. 

For real-world instruments, these timbral characteristics derive from the shape of the instrument, the physical properties of its material, and the way in which a musician plays the note~\cite{Benson_2006}. From a mathematical perspective, however, timbre can be understood by the phases and amplitudes of partials to the fundamental frequency. Partials are individual sine waves that make up the overall sound, and they can be harmonic (falling on or close to a member of the fundamental's harmonic series--also called overtones) or inharmonic, meaning that they do not align well with the harmonic series. Most pitched instruments are designed specifically so that their partials are harmonic, while percussive instruments and bells often have a more varied mix of both harmonic and inharmonic partials~\cite{maher_2008}. 

\begin{figure}[hbt!]
\centering
\begin{tikzpicture}

  \node[anchor=south west,inner sep=0] (image) at (0,0) 
       {\includegraphics[width=1.0\textwidth]{images/pianoflutegtrtimbre.png}};
  
  
  \begin{scope}[x={(image.south east)},y={(image.north west)}]
    
    \draw [black,dashed,-<] (0.136,0.02) -- (0.136,1.0) node[above] {440Hz};
  \end{scope}
\end{tikzpicture}
\caption{Graph of FFT Magnitudes showing the different timbres of piano, flute, and acoustic guitar playing a pitch of A440.}
    \label{fig:pianoflutegtrtimbre}
\end{figure}

\clearpage
\section{Frequency Modulation Synthesis}\label{sec:FM_synthesis}
In 1967, Stanford researcher John Chowning developed a new method of synthesizing audio using a technique called frequency modulation (FM) synthesis. In its most simple case, FM synthesis requires two oscillators, which each produce a tone at a given frequency in Hz. Instead of combining the oscillator outputs directly, we designate one as the ``carrier'' and one as the ``modulator'' and use the modulator's frequency $\w_m$ to modify the frequency parameter of the carrier $\w_c$ by an amount determined by the modulation index $\beta = B/\w_m$ where $B$ is the peak deviation from the original carrier. Mathematically, we can describe the instantaneous amplitude, or waveform, of the modulated carrier as:
\begin{equation}
  FM(t) = \sin(\w_ct + \beta \sin( \w_mt))  
\end{equation}
This small change to how the two oscillators interact has a tremendous impact on the character, or timbre, of the sound produced, being far richer and more complex than two oscillators arranged in parallel. The reason for this lies in the resulting spectrum of FM synthesis. Notably, higher modulation indices ``steal'' energy from the carrier. For some values, the carrier amplitude can be entirely reduced to zero. 

To better illustrate this, we can take the Fourier transform of this FM arrangement with an added carrier amplitude $A$ and deviation $B$. First however, we must integrate the instantaneous frequencies $\tau$ produced by the modulator over the duration of the signal $t$:
\begin{align*}
    FM(t) &= A\sin\bigg(\int_0^t (\w_c + B\sin(\w_m \tau))d\tau\bigg)\\
    &= A\sin\bigg( \w_ct - \frac{B}{\w_m}(\cos(\w_m t) - 1)\bigg)\\
    &= A\sin\bigg( \w_ct + \frac{B}{\w_m}(\sin(\w_m t - \pi/2) + 1)\bigg)
\end{align*}
Ignoring the constant phase terms on the carrier ($\varphi_c = \frac{B}{\w_m}$) and the modulator ($\varphi_m = -\pi/2$), we get the same equation as before~\cite{Chowning_1973}:
\begin{equation}
    FM(t) \approx A\sin(\w_c t + \beta \sin(\w_m t))
\end{equation}
In order to take the Fourier transform, note that functions of the form $\sin(\beta\sin(\theta))$ have the Fourier series~\cite{kreh_2012}:
\begin{equation}
    2\sum_{n=0}^\infty J_{2n+1}(\beta)\sin((2n+1)\theta)
\end{equation}
Where $J_n$ refers to the $n$-th Bessel function of the first kind~\cite{Benson_2006}.\footnote{See \textit{Music: A Mathematical Offering} Section 2.8-2.10 for more information and the appendix for useful properties.} More generally, we can expand $A\sin(\w_c t + \beta \sin(\w_m t))$ in terms of its Fourier series as:
\begin{equation}
    A\sin(\w_c t + \beta \sin(\w_m t)) = A\sum_{n=-\infty}^\infty J_n(\beta)\sin((\w_c + n\w_m)t)
\end{equation}
Notice that we have gone from each operator having a single nonzero term in its Fourier series to their composition now having infinitely many. Furthermore, we must note two important properties of Bessel functions that emerge from its power series definition found by applying the Frobenius method to Bessel's equation~\cite{abramowitz+stegun}:
\begin{equation}
    J_\alpha(x) = \sum_{m=0}^\infty\frac{(-1)^m}{m!\Gamma(m+\alpha+1)}\left(\frac{x}{2}\right)^{2m+\alpha}
\end{equation}
Where $\Gamma$ denotes the gamma function, or the generalized factorial function. In practice, we will only use the discrete case of $ \alpha \in \Z$ when describing frequency modulation, which is simply the factorial function. 

The first observation is that $J_n$ is even for even $n$ and odd for odd $n$, meaning that for even terms, our spectrum is symmetric about the carrier and for odd terms, the magnitudes of their Fourier coefficients are symmetric about the carrier although their signs are inverted. We call pairs of terms with summation indices $(-n,n) :n\neq 0$ ``sidebands.'' The second crucial observation is that for larger values of $\alpha$, and therefore $n$, the series converges to very small values because the gamma function in the denominator grows much faster than $x^\alpha$. What this means is that even for large modulation indices $\beta$, there will always be sidebands for which further series terms have a negligible impact on timbre. 
  
\chapter{Spectral Leakage and The Dirichlet Kernel}
\label{ch:3}
\section{Introduction}\label{sec:leakage}
One of the most relevant properties of the Fourier transform is that the operation of convolution in the Cartesian coordinate space is equivalent to multiplication after one has taken the Fourier transform:
\begin{equation}\label{eq:conv}
\mathcal{F}\{f \ast g\}(\omega) = \mathcal{F}\{f\}(\omega) \cdot \mathcal{F}\{g\}(\omega)
\end{equation}
The advantage of this is that it is often much easier and more computationally efficient to compute a Fourier transform of two functions than it is to perform their convolution~\cite{David13}.

However, when working in the discrete case, this property has consequences that do not always apply to the continuous Fourier transform. When dealing with a finite number of $N$ DFT samples, one implicitly applies a rectangular windowing function of length $N$ to the original function. This happens because trigonometric functions like sine and cosine extend infinitely, whereas a discrete function described by a finite number of samples does not.

Mathematically, when we compute the DFT of a signal $x[n]$, we are actually computing the DFT of $x[n] \cdot w[n]$, where $w[n]$ is a rectangular window function defined as:
\begin{equation}
w[n] = 
\begin{cases}
1, & 0 \leq n \leq N-1 \\
0, & \text{otherwise}
\end{cases}
\end{equation}
By the same convolution property, the DFT of the windowed signal is the convolution of the DFT of the original signal with the DFT of the window function~\cite{pereyra_ward_harmonic}:
\begin{equation}
\mathcal{F}\{x[n] \cdot w[n]\} = \mathcal{F}\{x[n]\} \ast \mathcal{F}\{w[n]\}
\end{equation}
The DFT of the rectangular window function $w[n]$ is the Dirichlet kernel $D_N[k]$, given by:
\begin{equation}
D_N[k]  = \frac{\sin(\pi k)}{\sin(\pi k/N)}
\end{equation}
Consequently, the frequencies $X[k] = \mathcal{F}\{x[n]\}$ get ``smeared out'' across the resulting spectrum after being convolved with $D_N[k] = \mathcal{F}\{w[n]\}$. This smearing effect can be precisely stated as follows: If $X[k]$ represents the true DFT of an infinite-length signal, and $\hat{X}[k]$ represents the DFT we actually compute from the windowed signal, then~\cite{9345707}:
\begin{equation}
\hat{X}[k] = \sum_{m=0}^{N-1} X[m] \cdot D_N[k-m]
\end{equation}
This equation shows that the value at each frequency bin $k$ in our computed DFT is not simply the energy at that frequency, but rather a weighted sum of the energy at all frequencies, where the weighting function is the Dirichlet kernel centered at $k$. This is the mathematical definition of spectral leakage—energy from one frequency ``leaks'' into surrounding frequency bins according to the shape of the Dirichlet kernel.

\begin{align*}
    \includegraphics[scale=0.3]{images/440curvefrontside.png}
    \includegraphics[scale=0.3]{images/440curvebackside.png}
\end{align*}

An interesting property of this spectral leakage is that for any given component, the observed leakage pattern depends only on the underlying sampling frequency and the true continuous parameters of the component itself. Another consequence of using a finite number of points in a DFT is that one must give up perfect frequency resolution; instead, nearby frequencies are summed into ``bins'' whose width $\Delta f$ is given by:
\begin{equation}
\Delta f = \frac{f_s}{N}
\end{equation}
where $f_s$ is the sample rate (samples per second) and $N$ is the number of samples.

To analyze the behavior of the Dirichlet kernel more precisely, define $\omega_k = \frac{2\pi k}{NT}$ where $k \in \R$ is the bin index extended to the real numbers. Here, $T = \frac{N}{f_s}$ is the total time interval of the signal. The spectral rectangular window (Dirichlet kernel) of an $N$-point DFT can be expressed as~\cite{Harris_1978}:
\begin{equation}
W(\omega) = \exp{\left(-i\frac{\omega T}{2}\right)}\frac{\sin{\left(\frac{N\omega T}{2}\right)}}{\sin{\left(\frac{\omega T}{2}\right)}}
\end{equation}

To understand how the kernel behaves with respect to frequency offsets, we can rewrite this in terms of the difference between a bin frequency $\omega_k$ and the true frequency $\omega_f$. Let $\omega_d = \omega_k - \omega_f = \frac{2\pi}{NT}(k - f)$ represent this offset. Substituting this into the kernel equation:
\begin{align}
W(\omega_d) &= \exp{\left(-i\frac{\omega_d T}{2}\right)}\frac{\sin{\left(\frac{N\omega_d T}{2}\right)}}{\sin{\left(\frac{\omega_d T}{2}\right)}} \\
\implies W_f(k)&= \exp{\left(-i\frac{\frac{2\pi}{NT}(k-f)T}{2}\right)}\frac{\sin{\left(\frac{N}{2}\frac{2\pi}{NT}(k-f)T\right)}}{\sin{\left(\frac{1}{2}\frac{2\pi}{NT}(k-f)T\right)}} \\
&= \exp{\left(-i\frac{\pi(k-f)}{N}\right)}\frac{\sin{\left(\pi(k-f)\right)}}{\sin{\left(\frac{\pi(k-f)}{N}\right)}}
\end{align}

This reveals an important property: when the difference between the bin frequency and true frequency is an integer (i.e. $k-f \in \mathbb{Z}$), the numerator term $\sin(\pi(k-f))$ becomes zero except when $k-f = 0$, resulting in:
\begin{equation}
k-f \in \mathbb{Z} \setminus \{0\} \implies \sin{\left(\pi (k-f)\right)} = 0 \implies W(\omega_d) = 0
\label{eq:coherent_kernel}
\end{equation}

This means that the frequency response of the Dirichlet kernel is zero at all integer bin offsets from the true frequency, which creates the characteristic pattern of nulls in the frequency response visible in Figures \ref{fig:stackexch} and \ref{fig:windows}~\cite{BREITENBACH1999135}. In practice, however, true frequencies rarely align exactly with the DFT bin frequencies, leading to spectral leakage.
\clearpage
\begin{figure}[hbt!]
  \centering
    \includegraphics[width=0.8\linewidth]{images/dsp_stack_exchange_IR.png}
  \caption{Frequency vs amplitude impulse response of the Dirichlet kernel~\cite{hotpaw_2013}}
  \label{fig:stackexch}
\end{figure}

In the worst case, when a true frequency falls exactly halfway between two bins, those bins exhibit equal response magnitudes. This configuration not only maximizes leakage~\cite{Harris_1978}, but also results in significantly reduced magnitude accuracy at these bins due to the steep drop-off of the kernel's frequency response. The magnitude at these bins will be approximately $2/\pi \approx 0.637$ times the true magnitude. Similar magnitude distortion occurs for all frequencies that do not align perfectly with the DFT bin centers.



\newpage
\section{Window Functions}

Given that spectral leakage can substantially hamper analysis in the frequency domain, other window functions than the rectangular window are commonly used for their reduced global leakage response at the cost of a wider main lobe. Some of the most common are the Hamming, Hann, Blackman, and Triangle (Fejer-Bartlet) window functions~\cite{Harris_1978}.
\begin{figure}[hbt!]
    \centering
    \includegraphics[width=1.05\linewidth]{images/WindowFunctions.png}   
    \includegraphics[width=1.05\linewidth]{images/Window Functions 2.png}
    \caption{Common window functions and their frequency responses \cite{Gjelstrup_2025}}
    \label{fig:windows}
\end{figure}
\clearpage
Although window functions are extremely helpful, they are still prone to certain fundamental limitations. Notice in Figure \ref{fig:windows} that although window functions reduce the spectral leakage, they do so at the cost of both frequency precision and amplitude accuracy in the central lobe and its adjacent bins. This is by far the largest drawback to using window functions, as for scenarios in which high frequency resolution is paramount, especially when dealing with low sample rates or short signals, the global improvements to spectral leakage might not justify the cost to local precision. 


\section{Related Approaches to Spectral Leakage Reduction}

While window functions are the most common approach to addressing spectral leakage, several alternative methods have been proposed that aim to preserve spectral resolution while minimizing leakage effects. This section reviews the most significant approaches and highlights their relative advantages and limitations.

\subsection{Traditional Windowing Techniques}

The most common approach to addressing spectral leakage is through windowing techniques as described by Jwo et al.~\cite{jwo2019windowing}. These techniques apply various window functions to the time-domain signal before the Fourier transform to reduce discontinuities at the signal boundaries. As noted by Jwo et al.: ``Windowing amplitude modulates the input signal so that the spectral leakage is evened out. Thus, windowing reduces the amplitude of the samples at the beginning and end of the window, altering leakage.'' Popular windowing functions include:

\begin{itemize}
    \item \textbf{Hanning windows}: Offer better side lobe suppression (-31.5 dB) at the cost of wider main lobes
    \item \textbf{Blackman windows}: Provide superior side lobe suppression (-57 dB) with even wider main lobes
    \item \textbf{Kaiser windows}: Offer parameterized trade-off between main lobe width and side lobe level
\end{itemize}

While effective for many applications, windowing techniques fundamentally sacrifice frequency resolution to reduce leakage. 

\subsection{Frequency Domain Phase Correction (FXT)}

Xu~\cite{xu2006algorithm} proposed the FXT algorithm that addresses spectral leakage by compensating for non-integer cycle counts in the analysis window. Unlike window functions that modify amplitude relationships, FXT applies phase corrections to ``twist'' the reference space, effectively making non-periodic signals appear periodic to the Fourier transform.

The approach works by:
\begin{enumerate}
    \item Identifying the integer and fractional components of signal cycles
    \item Demodulating by the integer component
    \item Calculating the phase difference caused by the fractional component
    \item Applying a twiddle function to compensate for this phase difference
\end{enumerate}

For multi-tone signals, FXT subdivides the spectrum and applies corrections independently to each subdivision. While effective for converter testing applications, this method does not explicitly extract or model individual spectral components and is primarily focused on correction rather than representation.

\subsection{Fundamental Identification and Replacement}

Wu et al.~\cite{wu2012new} developed an ``Improved Fundamental Identification and Replacement'' technique that focuses on spectral testing of analog-to-digital converters (ADC). Their approach:

\begin{enumerate}
    \item Identifies the parameters (amplitude, frequency, phase) of the fundamental sine component
    \item Replaces it with a coherent sine component having the same amplitude/phase but adjusted frequency
    \item Applies FFT to the modified signal with minimal leakage effects
\end{enumerate}

The method particularly excels for single-tone signals but, as the authors note, ``is not well qualified for multitone signal tests.'' Wu's approach includes optimizations for computational efficiency by selecting data lengths with only prime factors of 2 and 3.

\subsection{ESPRIT (Estimation of Signal Parameters via Rotational Invariance Techniques)}

ESPRIT~\cite{roy1989esprit} is a high-resolution parameter estimation technique that exploits the rotational invariance property of signal subspaces. Unlike Fourier-based methods that rely on windowed data, ESPRIT uses the eigenstructure of the signal's correlation matrix to identify signal components.

The key advantage of ESPRIT over traditional DFT-based approaches is its ability to estimate frequency parameters with sub-bin precision without suffering from spectral leakage. The method operates by:

\begin{enumerate}
    \item Constructing a data matrix from the time-domain signal samples
    \item Computing the signal subspace using eigendecomposition or singular value decomposition 
    \item Exploiting the rotational invariance property between subspaces to estimate signal parameters
    \item Solving for the frequencies, amplitudes, and phases of the underlying components
\end{enumerate}

While ESPRIT achieves remarkably high frequency resolution and effectively eliminates spectral leakage, it has several significant limitations. First, it assumes that the signal consists of a small number of complex sinusoids in white noise, which is often not the case for real-world signals with rich harmonic content. Second, ESPRIT's computational complexity scales cubically with the data size ($\mathcal{O}(n^3)$) due to its reliance on eigendecomposition, making it impractical for large datasets. Finally, it requires \textit{a priori} knowledge of the number of signal components, which is rarely available in practice.

\subsection{Periodogram Techniques}

As investigated by Jwo et al.~\cite{jwo2019windowing}, the Welch method can be used to reduce spectral leakage effects when combined with windowing. By dividing signals into overlapping segments, applying windows to each segment, and averaging the resulting periodograms, this approach achieves enhanced spectral resolution and reduces noise effects.

\subsection{Zero Padding to Minimize Kernel Effects}
Puche-Panadero et al.~\cite{9345707} present an approach to spectral leakage reduction that leverages zero padding to effectively minimize Dirichlet kernel effects. Their technique preserves the rectangular window's narrow main lobe advantages while still addressing leakage issues. The method works by creating a finer frequency grid through zero padding, which allows for strategic sampling of the spectrum at points where the Dirichlet kernel's sidelobes approach zero.

The approach involves padding the original signal with zeros to create a much finer ``fDFT'' (fine Discrete Fourier Transform) with frequency bins separated by $f'= f/R_N$, where $R_N$ is the padding ratio. From this fine spectrum, the authors then extract a ``Reduced Leakage DFT'' (RLDFT) by selecting bins at frequencies where the spectral leakage is minimized. This occurs specifically at frequencies that are offset from the fundamental component by integer multiples of the original bin width f. As demonstrated in their experimental validation on a 3.15 MW induction motor, this technique can reduce spectral leakage by more than 60 dB compared to conventional FFT approaches, making previously hidden fault harmonics clearly visible even at very low slip conditions.


\subsection{The Tabei-Ueda Method}
After developing the Dirichlet Kernel Deconvolution algorithm, the author discovered the Tabei-Ueda method~\cite{tabei1987high}, which addresses the related problem of high-precision parameter estimation from FFT data. Their approach uses a Hamming window before FFT computation and achieves sub-bin resolution through interpolation formulas derived from the window's frequency response. For simple sinusoidal signals, both methods produce nearly identical results, validating the mathematical soundness of our independently developed approach.
While the Tabei-Ueda method uses predetermined interpolation formulas for a Hamming-windowed signal, DKD takes a different approach by modeling spectral leakage as a convolution with the Dirichlet kernel and iteratively extracting components. This deconvolution perspective has enabled us to develop several novel signal processing techniques including the Frequency Modulation Transform (FMT), idealized resampling, and source separation methods that extend beyond the original parameter estimation problem.

\subsection{Limitations of Existing Approaches}

These methods, while addressing certain aspects of spectral leakage, have several limitations:
\begin{itemize}
    \item They often focus on correction rather than representation
    \item Many require \textit{a priori} knowledge about signal structure
    \item Most provide solutions specifically tailored to applications in a specific domain
    \item Few address the fundamental issue of precise sub-bin parameter identification
\end{itemize}

The approach presented in this thesis, Dirichlet Kernel Deconvolution (DKD), addresses these limitations by directly modeling the kernel response patterns and extracting components with sub-bin precision. Unlike previous methods that attempt to correct or mask leakage effects, DKD embraces the underlying mathematical structure of leakage to extract the true continuous-frequency components.

This enables not only leakage reduction but also applications in compression, resampling, and time-frequency analysis that other methods cannot readily support.
In the next chapter, we propose a novel windowless approach to addressing spectral leakage that obtains a ``best of both worlds'' solution where central-lobe fidelity is dramatically improved while also observing a near-complete reduction in spectral leakage. Given that applying the window before the transform is what causes leakage, we instead opt to analyze the signal after its FFT has been taken with a rectangular window. Then, through the utilization of properties of the Dirichlet kernel, we can obtain arbitrarily good approximations of the ``true'' (continuous) parameters of the components that make up the signal. The DKD algorithm is not intended to be perfectly optimized for performance in its current state, rather it should serve as theoretical foundations for a novel mathematical approach. 


\chapter{Dirichlet Kernel Deconvolution}
\label{ch:4}
Although the nature of frequency bins means that, in principle, we should not be able to determine the exact continuous frequency, phase, or amplitude for any given component looking only at the DFT, the spectral leakage at each bin actually depends on these hidden parameters. Thus, we can work backwards from the observed leakage in the DFT by exploiting the relationship between the underlying continuous components and their Dirichlet kernel responses to obtain the true parameters of each component. 

\section{Mathematical Foundations}
\label{subsec:mathematical_foundations}

The Dirichlet Kernel Deconvolution (DKD) method addresses a fundamental inverse problem in signal processing. Given a finite-length discrete Fourier transform $\mathbf{F} = \text{DFT}(\mathbf{S})$ of signal $\mathbf{S}$, we seek to decompose it into a set of components with parameter precision that transcends the resolution limitations of the standard DFT. In other words, we seek to reconstruct the continuous Fourier transform of the signal from its discrete counterpart.

Consider a real-valued signal composed of $M$ sinusoids. Its theoretical continuous Fourier series can be expressed as:
\begin{equation}
\label{eq:continuous_signal}
\mathbf{S}(t) = A_0 +\sum_{i=1}^{M} A_i \cos(2\pi f_i t + \varphi_i)
\end{equation}
where $A_i$, $f_i$, and $\varphi_i$ represent the amplitude, frequency, and phase of each component, respectively. When we sample this signal at a frequency of $f_s$ over $N$ points and compute its DFT, we obtain discrete frequency bins $\mathbf{F}_k$ that do not precisely match the true component frequencies due to both spectral leakage and interactions with binning.

However, this leakage follows a specific pattern determined by the Dirichlet kernel, which can be expressed for a specific frequency $f$ (in fractional bin units) and phase $\varphi$ as:
\begin{equation}
\label{eq:dirichlet_kernel}
D(f, \varphi)_k = \frac{\sin[\pi (f - k)]}{\sin[\pi (f - k)/N]} e^{-j(\pi (f - k) + \varphi)}
\end{equation}
where $k$ represents the bin index. For each component with parameters $(A_i, f_i, \varphi_i)$, the Dirichlet kernel generates a predictable pattern of energy distribution across multiple DFT bins.

The DKD algorithm approaches this as a sparse approximation problem, seeking to find the minimal set of components $\{(A_i^*, f_i^*, \varphi_i^*)\}_{i=1}^M$ such that:
\begin{equation}
\label{eq:sparse_approximation}
\mathbf{F} \approx \sum_{i=1}^{M} A_i^* D(f_i^*, \varphi_i^*)
\end{equation}
To solve this problem, we employ an iterative greedy approach similar to the CLEAN algorithm used in radio astronomy~\cite{Cornwell_2009}.\footnote{The CLEAN algorithm is the namesake of ``CleanFM.''}

We progressively identify and remove the component with the highest magnitude contribution to the residual spectrum.

\subsection{Spectral Properties of the Dirichlet Kernel}

To understand why our approach selects the bin with the highest magnitude, we must analyze the spectral properties of the Dirichlet kernel. The spectral window for the Dirichlet kernel under an $N$-point DFT is given by:
\begin{equation}
\label{eq:dirichlet_window}
W(\w_k) = \exp\left(-i\frac{N-1}{2}\w_k\right)\frac{\sin\left[\frac{N}{2}\w_k\right]}{\sin\left[\frac{1}{2}\w_k\right]}
\end{equation}
In terms of fractional bins where $\w_k = \frac{2\pi}{NT}k$, we can rewrite this as:
\begin{align}
\label{eq:dirichlet_window_fractional}
W(kT) &= \exp\left(-i\frac{N-1}{2}\frac{2\pi}{NT}kT\right)\frac{\sin\left[\frac{N}{2}\frac{2\pi}{NT}kT\right]}{\sin\left[\frac{1}{2}\frac{2\pi}{NT}kT\right]}\\
&= \exp\left(-i\left[\frac{N}{2}\frac{2\pi}{NT}kT + \frac{1}{2}\frac{2\pi}{NT}kT\right]\right)\frac{\sin\left[\frac{N}{2}\frac{2\pi}{NT}kT\right]}{\sin\left[\frac{1}{2}\frac{2\pi}{NT}kT\right]} \\
&= \exp\left(-i\left[\pi k + \frac{\pi k}{N}\right]\right)\frac{\sin\left[\pi k\right]}{\sin\left[\frac{\pi k}{N}\right]}
\end{align}

By Euler's formula, we can express this in terms of its real and imaginary components:
\begin{equation}
\label{eq:dirichlet_complex}
W(kT) = \frac{\sin\left[\pi k\right]}{\sin\left[\frac{\pi k}{N}\right]}\left(\cos\left[\pi k + \frac{\pi k}{N}\right] - i\sin\left[\pi k + \frac{\pi k}{N}\right]\right)
\end{equation}
The magnitude of this complex response is given by:
\begin{align}
\label{eq:dirichlet_magnitude}
\|W(kT)\|_2 &= \left|\frac{\sin\left[\pi k\right]}{\sin\left[\frac{\pi k}{N}\right]}\right| \sqrt{\cos^2\left[\pi k + \frac{\pi k}{N}\right] + \sin^2\left[\pi k + \frac{\pi k}{N}\right]} \\
&= \left|\frac{\sin\left[\pi k\right]}{\sin\left[\frac{\pi k}{N}\right]}\right|
\end{align}
where we have applied the Pythagorean identity in the final step.

For a true frequency that aligns exactly with a bin ($k = 0$), we have:
\begin{equation}
\label{eq:dirichlet_at_zero}
\lim_{k\to 0} \left|\frac{\sin\left[\pi k\right]}{\sin\left[\frac{\pi k}{N}\right]}\right| = N
\end{equation}
Furthermore, as $k$ increases, the magnitude of the Dirichlet kernel decreases. Specifically:
\begin{equation}
\label{eq:dirichlet_bounds}
0 \leq \left|\frac{\sin\left[\pi k\right]}{\sin\left[\frac{\pi k}{N}\right]}\right| \leq \frac{1}{\left|\sin\left[\frac{\pi k}{N}\right]\right|}
\end{equation}
Since $|\sin x|$ is strictly increasing on the interval $(0, \frac{\pi}{2})$, we know that $\frac{1}{\left|\sin\left[\frac{\pi k}{N}\right]\right|}$ must be decreasing on that same interval.

For each bin $m$ away from the true frequency, the fractional bin offset increases consistently by 1. Given that $|\sin[\pi(m + k)]| = |\sin[\pi m + \pi k]| = |\sin[\pi k]|$ for all $m \in \mathbb{Z}$, the amplitude at each subsequent bin is strictly decreasing as we move away from the true frequency.

The one caveat is if the true frequency falls exactly between two bins, in which case the Dirichlet kernel response will be identical for these two bins. This, however, is not an issue for the algorithm, as the true frequency will still be found by the correlation search on the interval $[i-0.5, i+0.5]$ as long as we include the boundaries.
\clearpage
\subsection{Component Extraction Process}

For each iteration, we identify the bin $i$ with maximum magnitude in the current residual:
\begin{equation}
\label{eq:peak_bin}
i = \underset{\text{bin} \in [1, N/2-1]}{\mathrm{argmax}}\ \|\mathbf{R}_{\text{bin}}\|_2 = \underset{\text{bin}}{\mathrm{argmax}}\ \sqrt{\Re(\mathbf{R}_{\text{bin}})^2 + \Im(\mathbf{R}_{\text{bin}})^2}
\end{equation}

Based on the aforementioned properties of the Dirichlet kernel, we know that the true frequency of the dominant component will be within $\pm 0.5$ bins (or $\pm f_s/(2N)$ Hz) of the selected bin $i$.

To determine the precise frequency, we define a candidate frequency $f_i' = i + \varepsilon_i$, where $\varepsilon_i \in [-0.5, 0.5]$ is a fractional bin offset. The optimal offset $\varepsilon_i^*$ is found by maximizing the normalized cross-correlation:

\begin{equation}
\label{eq:optimal_freq}
\varepsilon_i^* = \argmax_{\varepsilon_i \in [-0.5, 0.5]} \left\{ \frac{\left|\sum_{k=i-\ell_f}^{i+\ell_f} \mathbf{R}_k \cdot \overline{D(i + \varepsilon_i, \varphi_i)_k}\right|}{\|\mathbf{R}_{[i-\ell_f, i+\ell_f]}\|_2 \cdot \|D(i + \varepsilon_i, \varphi_i)_{[i-\ell_f, i+\ell_f]}\|_2} \right\}
\end{equation}

where $\ell_f$ defines a limited correlation window around bin $i$ and $\varphi_i$ is the initial phase estimate.

Similarly, the optimal phase $\varphi_i^*$ is determined by:

\begin{equation}
\label{eq:optimal_phase}
\varphi_i^* = \argmax_{\varphi \in [-\pi, \pi]} \left\{ \frac{\Re\left(\sum_{k=i-\ell_\varphi}^{i+\ell_\varphi} \mathbf{R}_k \cdot \overline{D(i + \varepsilon_i^*, \varphi)_k}\right)}{\|\mathbf{R}_{[i-\ell_\varphi, i+\ell_\varphi]}\|_2 \cdot \|D(i + \varepsilon_i^*, \varphi)_{[i-\ell_\varphi, i+\ell_\varphi]}\|_2} \right\}
\end{equation}

Finally, the true amplitude is calculated by applying a correction factor that accounts for spectral leakage:
\begin{equation}
\label{eq:amplitude_correction}
A_i^* = A_i \cdot \frac{\sin[\pi \varepsilon_i^*/N]}{\sin[\pi \varepsilon_i^*]}
\end{equation}

where $A_i$ is the observed amplitude at bin $i$. This process is repeated iteratively until a termination condition is met, resulting in a set of components that accurately represent the continuous frequencies comprising the original signal.
\clearpage
\section{DKD: Single-Threaded}


\subsection{Algorithm Implementation}
\label{subsec:algorithm_implementation}

The implementation of the DKD algorithm builds directly upon the mathematical foundations established in Section~\ref{subsec:mathematical_foundations}. Here, we detail the complete control flow and computational procedure of the algorithm.

\subsubsection{Initialization and Preprocessing}

The control flow for Dirichlet Kernel Deconvolution begins by reading an input waveform file and computing its FFT using a custom implementation of the \texttt{dr\_wav} library~\cite{dr_wav}. From this process, we extract the FFT size $N$ and sample rate $f_s$, which are automatically determined from the file metadata.

Before the main decomposition routine begins, we perform several preprocessing steps:
\begin{enumerate}
    \item Copy the first $N/2$ bins from the FFT vector into a new residual vector $\mathbf{R}$.
    \item Set the first index in the residual to zero: $\mathbf{R}_0 = 0$.
\end{enumerate}
We restrict the analysis to the first $N/2$ bins because the DFT has Hermitian symmetry for real-valued input data~\cite{xu2006algorithm}. In other words, the spectrum has conjugate symmetry about the Nyquist bin—the FFT bin corresponding to the signal's Nyquist frequency. According to the Nyquist-Shannon Sampling Theorem (\ref{sec:NQST}), a signal cannot contain frequency components above half the sampling rate if it is to be perfectly reconstructed. For the signals analyzed in this study (typically sampled at 48 kHz or 44.1 kHz), the Nyquist bin represents frequency content at exactly $f_s/2$.

The information contained in the latter half of the FFT is essentially redundant due to this symmetry property, with a few exceptions related to mirrored effects across the Nyquist frequency that will be addressed later in \ref{subsubsec:cross_nyquist_kernel_reflections}. By operating only on the first $N/2$ bins, we substantially improve algorithmic efficiency, noting that we must take the spectrum's conjugate during reconstruction.

We set the DC component ($\mathbf{R}_0$) to zero because it corresponds to the average value of the signal and can be handled separately from the spectral decomposition process.

\subsubsection{Main Decomposition Loop}

After initializing the residual vector, we proceed with the main iterative loop:

\begin{enumerate}
    \item Identify the bin index $i$ with the highest magnitude in the current residual:
    \begin{equation}
    i = \argmax_{\text{bin} \in [1, N/2-1]} \|\mathbf{R}_{\text{bin}}\|_2
    \end{equation}
    \item Check termination conditions: if the maximum magnitude falls below a threshold $\epsilon$ or the maximum number of components $M_{\max}$ is reached, exit the loop.
    \item For the identified peak bin $i$, perform frequency optimization to find the precise fractional bin frequency $f_i^* = i + \varepsilon_i^*$ that maximizes the cross-correlation between the Dirichlet kernel and the residual.
    \item Using the optimized frequency $f_i^*$, perform phase optimization to find $\varphi_i^*$.
    \item Calculate the true amplitude $A_i^*$ using the correction factor from Equation~\ref{eq:amplitude_correction}.
    \item Store the extracted component parameters $(A_i^*, f_i^*, \varphi_i^*)$ for later reconstruction.
    \item Update the residual by subtracting the Dirichlet kernel response of the extracted component:
    \begin{equation}
    \mathbf{R}_k = \mathbf{R}_k - A_i^* \cdot D(f_i^*, \varphi_i^*)_k \quad \forall k \in [1, N/2-1]
    \end{equation}
    \item Return to step 1 and repeat until termination conditions are met.
\end{enumerate}
\begin{figure}[hbt!]
    \centering
    \includegraphics[width=0.8\linewidth]{images/DKD_ST_CF.png}
    \caption{Control flow diagram for the main DKD decomposition loop}
    \label{fig:control_flow}
\end{figure}

This iterative approach progressively refines the residual, extracting components in order of their energy contribution to the signal.
\clearpage

\subsection{Golden Section Search}
\label{subsubsec:golden_section_search}

Rather than evaluating all possible fractional bin offsets, we employ golden section search~\cite{Kiefer_1953} to efficiently find the optimal frequency and phase. The golden section search is particularly well-suited for this problem for several reasons:
\begin{itemize}
    \item For sparse spectra, the cross-correlation function has a single maximum in the interval $[-0.5, 0.5]$ for frequency optimization, which similarly holds for phase.\footnote{The presence of multiple maxima is extremely rare, even in dense signals; however, DKD is actually robust to errors associated with subtracting an incorrect component from the residual. Although the parameter estimation for that component will be inaccurate, the algorithm will still converge to the target spectrum. This is discussed further in Section \ref{sec:dkdni}. }
    \item The function is relatively smooth but lacks a simple analytical derivative.
    \item The method converges logarithmically, offering substantial computational savings over grid search and allows for adjustable termination thresholds depending on the application.
\end{itemize}

The golden section search algorithm exploits the golden ratio $\varphi = \frac{1+\sqrt{5}}{2} \approx 1.618$ to optimally partition the search interval:

\begin{enumerate}
    \item Initialize search interval $[a, b] = [i-0.5, i+0.5]$ for frequency optimization.
    \item Calculate interior points $c = b - (b - a)/\varphi$ and $d = a + (b - a)/\varphi$.
    \item Evaluate the objective function (normalized cross-correlation) at points $c$ and $d$.
    \item If $f(c) > f(d)$, set $b = d$; otherwise, set $a = c$.
    \item Repeat until the interval width $|b - a|$ is below the tolerance threshold $\epsilon_{\text{GSS}}$.
\end{enumerate}

This approach reduces the complexity from $O(N \cdot K)$ to $O(N \cdot \log K)$, where $K$ is the number of candidate values that would be evaluated in a grid search.
\clearpage
\subsection{Determining Optimal Parameters}
\label{subsubsec:determining_optimal_parameters}
The performance of the DKD algorithm depends on several key parameters:
\begin{enumerate}
    \item Frequency correlation window size ($\ell_f$): Controls the number of bins considered when optimizing frequency.
    \item Phase correlation window size ($\ell_\varphi$): Controls the number of bins considered when optimizing phase.
    \item Golden section search tolerance ($\epsilon_{\text{GSS}}$): Determines the precision of frequency and phase optimization.
    \item Maximum component count ($M_{\max}$): Limits the total number of components extracted.
    \item Residual threshold ($\epsilon_{\text{res}}$): Determines when to terminate the extraction process based on residual energy.
\end{enumerate}
Through empirical testing, we have determined optimal values for these parameters:
\begin{enumerate}
    \item $\ell_f = 4$: This provides sufficient context to accurately determine frequency while limiting computational overhead.
    \item $\ell_\varphi = 3$: A slightly smaller window is sufficient for phase optimization since frequency has already been determined.
    \item $\epsilon_{\text{GSS}} = 10^{-8}$: This provides sub-millihertz frequency precision for typical sample rates.
    \item $M_{\max} = 20000$: This upper bound is usage specific, but for a complex, multi-source audio signal of 3s at 48kHz, 20K corresponds to a 99.9\% total energy capture. 
    \item $\epsilon_{\text{res}} = 10^{-6} \cdot E_{\text{total}}$: The algorithm terminates when the residual energy falls below one-millionth of the initial total energy. (Determined experimentally by running DKD with $M_{\max}=1$ on signals composed of a single sine wave)
\end{enumerate}
These values represent a balance between accuracy and computational efficiency. For signals with very closely spaced frequency components, larger correlation windows may be beneficial, while simpler signals may permit different windows without loss of accuracy. 

\subsection{Cross-Nyquist Kernel Reflections}
\label{subsubsec:cross_nyquist_kernel_reflections}

A critical aspect of the DKD algorithm is proper handling of frequencies near the Nyquist boundary. When a component has frequency content that approaches the Nyquist limit, its spectral leakage pattern reflects across it, creating additional mirrored contributions to the spectrum.

For each component with parameters $(A_i^*, f_i^*, \varphi_i^*)$, we must consider both its direct contribution to the spectrum and its reflection. The reflection appears at bin $N-k$ for each bin $k$ and has conjugate-symmetric properties. 
During the residual update step, we might naively handle these reflections:
\begin{equation}
\mathbf{R}_{N-k} = \overline{\mathbf{R}_k} \quad \forall k \in [1, N/2-1]
\end{equation}
\begin{figure}[hbt!]
    \centering
    \includegraphics[width=0.8\linewidth]{images/cross_nyquist_final.png}
    \caption{Comparison of continuous DKD spectrum without accounting for Cross-Nyquist kernel reflections before the IFFT (red) and ground-truth binned 10kHz signal (blue).}
    \label{fig:cross_nyquist}
\end{figure}
However, this does not capture the kernel effects from the Nyquist-reflected $N-k$-th bin back to the $k$-th bin. To account for this, we must instead compute:
\begin{equation}
\mathbf{R}_{k} =\mathbf{R}_{k} + \overline{A_i^*\cdot D(f^{*'}_i, \varphi^{*'}_i)_{N-k}} \quad \forall k \in [0, N-1]
\end{equation}
Where $f^{*'}_i = f^{*}_i - N +m$ and $\varphi_i^{*'} + \pi f^{*'}_i / N$. This also ensures that the conjugate symmetry property of the DFT is maintained, which is essential for the reconstructed signal to remain real-valued. Without proper Cross-Nyquist handling, components, especially those near the Nyquist frequency, would be incorrectly estimated, leading to artifacts in both the decomposition and reconstruction.\footnote{Interestingly, Cross-Nyquist reflections were also discovered independently by Tabei and Udea and briefly mentioned in their paper~\cite{tabei1987high}, but the Hamming window prevented most of these reflections except when dealing with frequencies near the Nyquist bin (or with extremely low values of $N$). However, given the Dirichlet kernel affects the entire frequency domain, DKD must address this behavior for every component.}
\begin{figure}[hbt!]
    \centering
    \includegraphics[width=0.8\linewidth]{images/cross_nyquist_after2.png}
    \caption{Comparison of continuous DKD spectrum with Cross-Nyquist kernel reflections before the IFFT (red, covered) and ground-truth binned 10kHz signal (blue).}
    \label{fig:cross_nyquist_after}
\end{figure}
\newpage
\subsection{Reconstruction}
\label{subsubsec:reconstruction}

Once we have extracted a set of $M$ components $\{(A_i^*, f_i^*, \varphi_i^*)\}_{i=1}^M$, we can reconstruct the original signal in either the frequency or time domain.

For frequency-domain reconstruction, we generate a synthetic spectrum by summing the Dirichlet kernel responses for all components:
\begin{equation}
\mathbf{F}_{\text{recon},k} = \sum_{i=1}^{M} A_i^* \cdot D(f_i^*, \varphi_i^*)_k \quad \forall k \in [0, N-1]
\end{equation}
Ensuring conjugate symmetry across the Nyquist boundary and applying Cross-Nyquist reflections:
\begin{equation}
\mathbf{F}_{k} =\mathbf{F}_{k} + \overline{A_i^*\cdot D(f^{*'}_i, \varphi^{*'}_i)_{N-k}} \quad \forall k \in [0, N-1]
\end{equation}
The reconstructed time-domain signal is then obtained by applying the inverse FFT:
\begin{equation}
\mathbf{S}_{\text{recon}}(t) = \text{IFFT}(\mathbf{F}_{\text{recon}})
\end{equation}
Alternatively, for time-domain reconstruction, we can directly synthesize the signal using the extracted parameters:
\begin{equation}
\mathbf{S}_{\text{recon}}(t) = \sum_{i=1}^{M} A_i^* \sin(2\pi \frac{f_i^* f_s}{N} t + \varphi_i^*)
\end{equation}
where $f_s$ is the sampling frequency and $\frac{f_i^* f_s}{N}$ converts the fractional bin frequency to Hz.

The reconstructed signal provides a cleaned version of the original, with only the most significant components retained. This can be valuable for applications such as noise reduction, compression, and signal analysis.

\subsection{Complexity Analysis}

The time complexity of the single-threaded DKD algorithm is $O(M \cdot N \cdot \ell \cdot \log K)$, where:
\begin{itemize}
    \item $M$ is the number of components extracted
    \item $N$ is the FFT size
    \item $\ell$ is the correlation window size
    \item $K$ is the precision of the frequency/phase search (determined by $\epsilon_{\text{GSS}}$)
\end{itemize}

The space complexity is $O(N + M)$, which accounts for the residual spectrum and the storage of extracted components.

This represents a significant improvement over naive approaches that would require $O(M \cdot N^2 \cdot K)$ operations. The key optimizations that enable this performance are:
\begin{enumerate}
    \item Limited correlation windows ($\ell \ll N$)
    \item Golden section search ($\log K$ instead of $K$ evaluations)
    \item Working with only the first $N/2$ bins due to conjugate symmetry
\end{enumerate}

In practice, the algorithm's performance is also influenced by implementation details such as cache efficiency, vectorization, and memory access patterns. These considerations become increasingly important for large FFT sizes. 

\section{DKD: Multiple Threads}

While the single-threaded DKD algorithm provides excellent spectral decomposition, its computational complexity becomes a limiting factor for large signals. In this section, we present a parallel extension of the DKD algorithm that leverages multi-core processing to significantly reduce execution time while maintaining decomposition quality.

\subsection{Control Flow}

The parallel DKD algorithm introduces several key modifications to the single-threaded approach, specifically designed to enable concurrent component extraction with minimal synchronization overhead.

\subsubsection{Parallel Processing Strategy}

Our parallelization strategy centers around the observation that frequency components sufficiently separated in the spectrum can be extracted concurrently with minimal interference. We define a parallel workflow as follows:

\begin{enumerate}
    \item \textbf{Initialization:} As in the single-threaded version, we prepare the residual spectrum from the input signal's FFT.
    
    \item \textbf{Frequency Band Division:} Divide the spectrum into multiple bands, with each band assigned to a separate thread. We partition the frequency range $[1, N/2-1]$ into $B$ bands where $B$ is determined by the available hardware threads.
    
    \item \textbf{Concurrent Peak Identification:} Each thread independently locates local peak magnitudes within its assigned band.
    
    \item \textbf{Batch Component Selection:} Collect peaks from all bands and select a subset of peaks that are spectrally well-separated.
    
    \item \textbf{Parallel Component Optimization:} Optimize frequency, phase, and amplitude for all selected components concurrently.
    
    \item \textbf{Synchronized Residual Update:} Apply batch updates to the shared residual spectrum.
    
    \item \textbf{Iteration:} Repeat steps 3-6 until convergence criteria are met.
\end{enumerate}

This approach allows us to extract multiple components per iteration, significantly accelerating the overall decomposition process.
\newpage % Might remove newpage if it looks bad after modifying
\subsubsection{Thread Pool Implementation}

To manage the parallel execution, we implement a custom thread pool as shown in the following excerpt:

\begin{lstlisting}[language=C++, 
                   caption=Thread Pool Implementation, 
                   basicstyle=\footnotesize\ttfamily,
                   breaklines=true,
                   breakatwhitespace=true,
                   postbreak=\mbox{\textcolor{red}{$\hookrightarrow$}\space}]
class ThreadPool {
private:
    std::vector<std::thread> workers;
    std::queue<std::function<void()>> tasks;
    std::mutex queue_mutex;
    std::condition_variable condition;
    bool stop;

public:
    ThreadPool(size_t threads) : stop(false) {
        for (size_t i = 0; i < threads; ++i) {
            workers.emplace_back([this] {
                while (true) {
                    std::function<void()> task;
                    {
                        std::unique_lock<std::mutex> lock(queue_mutex);
                        condition.wait(lock, [this] { 
                            return stop || !tasks.empty(); 
                        });
                        if (stop && tasks.empty()) return;
                        task = std::move(tasks.front());
                        tasks.pop();
                    }
                    task();
                }
            });
        }
    }

    template<class F>
    void enqueue(F&& f) {
        {
            std::unique_lock<std::mutex> lock(queue_mutex);
            tasks.emplace(std::forward<F>(f));
        }
        condition.notify_one();
    }

    ~ThreadPool() {
        {
            std::unique_lock<std::mutex> lock(queue_mutex);
            stop = true;
        }
        condition.notify_all();
        for (auto &worker : workers) {
            worker.join();
        }
    }
};
\end{lstlisting}

The thread pool manages a fixed number of worker threads that process tasks from a shared queue. This design allows us to efficiently distribute workloads across available cores while minimizing thread creation/destruction overhead.

\subsubsection{Batch Processing}

A key innovation in our parallel implementation is the batch processing of components. Rather than extracting components strictly sequentially by magnitude, we extract multiple well-separated components in each iteration:
\begin{enumerate}
    \item Identify multiple peak candidates across the spectrum.
    \item Filter candidates using spectral distance constraints (explained in Section \ref{subsec:spectral_distance}).
    \item Optimize and extract the filtered components concurrently.
    \item Update the residual with all extracted components.
\end{enumerate}
The batch size is chosen adaptively based on the remaining spectral energy and available processing units, with a maximum value determined by \texttt{BATCH\_SIZE}.

\subsection{Spectral Distance}
\label{subsec:spectral_distance}

A critical aspect of parallel DKD is ensuring that simultaneously extracted components do not significantly interfere with each other. We define a spectral distance metric that quantifies the separation between peaks:
\begin{equation}
d(p_1, p_2) = \frac{|f_1 - f_2|}{N}
\end{equation}
where $f_1$ and $f_2$ are the fractional bin frequencies of peaks $p_1$ and $p_2$, and $N$ is the FFT size. This normalized distance represents the relative separation in the frequency domain.

For parallel extraction, we require a minimum spectral distance $d_{\min}$ between any two simultaneously processed peaks:
\begin{equation}
d_{\min} = \frac{2f_s}{N}
\end{equation}
where $f_s$ is the sample rate. This threshold ensures that the Dirichlet kernel responses of concurrently processed components have minimal overlap in their most significant regions. Specifically, we want to avoid processing two peaks with a shared adjacent bin.

When selecting peaks for batch processing, we apply a greedy approach:
\begin{enumerate}
    \item Sort all candidate peaks by magnitude (descending).
    \item Select the highest magnitude peak.
    \item Add subsequent peaks only if they maintain the minimum spectral distance from all previously selected peaks.
    \item Continue until reaching the maximum batch size or exhausting candidates.
\end{enumerate}
This selective process ensures that we can safely optimize and extract multiple components without introducing significant artifacts due to interference between their Dirichlet kernel responses.
\subsection{Parallel Optimization}

Each worker thread in the pool independently optimizes component parameters for its assigned peaks using the same golden section search and correlation-based approach described in Section \ref{subsubsec:golden_section_search}:

\begin{lstlisting}[language=C++, 
                   caption=Peak Optimization Worker Function, 
                   basicstyle=\footnotesize\ttfamily,
                   breaklines=true,
                   breakatwhitespace=true,
                   postbreak=\mbox{\textcolor{red}{$\hookrightarrow$}\space}]
void optimizePeak(
    const std::vector<Complex>& residual,
    std::vector<CleanDFT::Component>& batch_components,
    const std::vector<PeakInfo>& selected_peaks,
    size_t peak_index,
    size_t N)
{
    // Get the peak information
    const PeakInfo& peak = selected_peaks[peak_index];
    // Find optimal frequency
    double true_freq = CleanDFT::findOptimalFrequency(
        residual, peak.bin, peak.bin);
    // Find optimal phase
    double true_phase = CleanDFT::findOptimalPhase(
        residual, peak.bin, true_freq);
    // Compute amplitude
    double amplitude = DirichletKernel::getAmplitudeAtBin(
        true_freq, std::abs(residual[peak.bin]), N, peak.bin);
    // Store result (no mutex needed - writing to separate indices)
    batch_components[peak_index] = {true_freq, true_phase, amplitude};
}
\end{lstlisting}

This function runs independently for each peak, allowing all peaks in a batch to be optimized concurrently across multiple threads.
\newpage
\subsection{Synchronized Residual Update}

After optimizing all components in a batch, we must update the shared residual spectrum. This update requires synchronization to avoid data races. We divide this task into two steps:

\begin{enumerate}
    \item \textbf{Independent Updates:} Each thread updates a specific portion of the residual for all extracted components in the batch.
    
\begin{lstlisting}[language=C++, 
                   caption=Parallel Residual Update Function, 
                   basicstyle=\footnotesize\ttfamily,
                   breaklines=true,
                   breakatwhitespace=true,
                   postbreak=\mbox{\textcolor{red}{$\hookrightarrow$}\space}]
void updateResidualRange(
    std::vector<Complex>& residual,
    const CleanDFT::Component& component,
    size_t start_bin,
    size_t end_bin,
    size_t N)
{
    for (size_t i = start_bin; i < end_bin; i++) {
        // Calculate frequency difference
        double diff = component.true_frequency - static_cast<double>(i);
        
        // Calculate phase correction
        double corrected_phase = component.true_phase 
                               + M_PI * diff / static_cast<double>(N);
        
        // Subtract component response from residual
        residual[i] -= component.amplitude * 
            DirichletKernel::getValueAtBin(diff, corrected_phase, N);

        // Handle mirror frequency (if needed)
        size_t mirror_bin = N - i;
        if (mirror_bin < N && mirror_bin > i) {
            double mirror_diff = component.true_frequency 
                               - static_cast<double>(mirror_bin);
            double mirror_phase = component.true_phase 
                                + M_PI * mirror_diff / static_cast<double>(N);
            
            // Mirror uses complex conjugate
            residual[mirror_bin] -= std::conj(component.amplitude *
                DirichletKernel::getValueAtBin(mirror_diff, mirror_phase, N));
        }
    }
}
\end{lstlisting}
    \clearpage
    \item \textbf{Cross-Nyquist Updates:} After all ranges are processed, we maintain conjugate symmetry by enforcing:
    \begin{equation}
    \mathbf{R}_{k} =\mathbf{R}_{k} + \overline{A_i^*\cdot D(f^{*'}_i, \varphi^{*'}_i)_{N-k}} \quad \forall k \in [0, N-1]
    \end{equation}
\end{enumerate}

By partitioning the residual update across frequency ranges, we minimize contention between threads while ensuring correct handling of all frequency components.

\subsection{Parameter Evaluation}

The performance of parallel DKD depends on several critical parameters:

\begin{enumerate}
    \item \textbf{Number of Threads:} The optimal thread count typically matches the number of physical CPU cores. Exceeding this can lead to diminishing returns due to context switching overhead.
    
    \item \textbf{Batch Size:} This controls how many components are extracted per iteration. Larger batches increase parallelism but may reduce accuracy if too many components interfere.
    
    \item \textbf{Spectral Distance Threshold:} Higher thresholds enforce greater separation between concurrent components, improving accuracy at the cost of parallel efficiency.
    
    \item \textbf{Frequency Bands:} The number of frequency bands for initial peak detection affects load balancing and peak identification accuracy.
\end{enumerate}

Through empirical testing across various signal types and hardware configurations, we determined the following optimal values:

\begin{table}[h]
\centering
\begin{tabular}{|l|c|c|}
\hline
\textbf{Parameter} & \textbf{Default Value} & \textbf{Recommended Range} \\
\hline
Batch Size & 256 & 128-1024 \\
\hline
Minimum Spectral Distance & $2f_s/N$ & $(0.5-2.0) \times f_s/N$ \\
\hline
Number of Bands & $\text{num\_threads}$ & $[1, 2] \times \text{num\_threads}$ \\
\hline
\end{tabular}
\caption{Optimal parameter values for parallel DKD implementation}
\end{table}

These values offer a good balance between performance and accuracy across a wide range of signal characteristics.

\subsection{Scalability Analysis}

To evaluate the scalability of our parallel implementation, we analyze both theoretical and empirical performance metrics:
\subsubsection{Practical Scaling}

In practice, our implementation achieves near-linear scaling up to the number of physical CPU cores. Empirically, we observe that running the the algorithm on 12 Apple M3 Max performance cores\footnote{Using the full 16 cores (12 performance and 4 efficiency cores) actually yields worse performance than simply using the 12 performance on their own. This is likely because the efficiency cores have a notably lower clock speed (2.8GHz vs 4.06GHz). This behavior is machine specific, and could be circumvented by using a system with identical clock speeds on all cores.} gives between 8-10$\times$ speedup at a negligible cost to accuracy (especially at high component counts and for long audio files). However, with such a substantial runtime improvement, this is well worth the minor tradeoff. One possible drawback is that there is a chance that batching/multithreading ``overshoots'' the true number of frequencies in the signal if it is sufficiently sparse; however, for most real-world signals, this is unlikely do to noise/interference.

Several factors limit perfect scaling:
\begin{enumerate}
    \item Synchronization overhead during residual updates
    \item Compounding losses to norm capture as batch size increases
    \item Load imbalance when processing non-uniform spectral distributions
    \item Memory bandwidth limitations for very large FFT sizes
    \item Batch size constraints due to spectral distance requirements
\end{enumerate}

Despite these limitations, the parallel DKD implementation delivers substantial performance improvements that make high-quality spectral decomposition practical for complex signals in real-world applications.

\section{Results}

The effectiveness of both single-threaded and multi-threaded DKD implementations can be evaluated across three key dimensions: performance, accuracy, and scalability. In this section, we present comprehensive results from our experimental evaluation.
\subsection{Overview and Qualitative Assessment}
The DKD algorithm accomplishes its goal of achieving a highly accurate approximation to signals of varying lengths, complexities, and resolutions by decomposing its input into a sum of its individual kernels. The algorithm in its current form is surprisingly fast despite minimal optimizations beyond multithreading, and can achieve high accuracy on both objective and perceptual metrics with component count proportional to signal complexity and inversely proportional to sample rate.  
\subsubsection{Single Frequencies}
For sufficiently sparse signals ($<$500 frequencies), the accuracy of DKD makes it an appealing compression/encoding scheme in cases when highly precise true frequencies are required. In the base case, running DKD on a 440Hz signal generated by an oscillator in Logic Pro X sampled at 48kHz with 441975 samples achieves a normalized root mean squared error of 1.47e-07. 
\begin{figure}[hbt!] 
\centering
    \includegraphics[width=0.7825\linewidth]{images/recon_single_zoomed.png}    
    \caption{Reconstructed DFT (yellow) completely covers the original (purple) }
\end{figure}
\clearpage

Recall that what we are essentially doing is fitting the image of the Dirichlet kernel in $\R\times\C$ to the observed set of points in $\C^N$. To better visualize what the optimal kernel parameters look like, we can plot its graph in three dimensions:
\begin{figure}[hbt!] 
\centering
    \includegraphics[width=0.96\linewidth]{images/440curvefrontside.png}    
    \caption{Complex graph of DKD with one component where $f^* = 440$Hz}
    \label{fig:complexir}
\end{figure}
\newline
As we can see in Figure \ref{fig:complexir}, the continuous kernel (red) only depends on the frequency, phase, and amplitude of the true continuous component. The measured FFT points (blue) will fall at points on the kernel corresponding to Equation \ref{eq:dirichlet_kernel}. However, as mentioned in \ref{eq:coherent_kernel}, if the frequency is an integer multiple of the sample frequency, then the auxiliary points collapse to zero. This will happen for a 1-second 480Hz signal sampled at 48kHz ($N = 48000$), which we can see in Figure \ref{fig:complexcoherent}.
\begin{figure}[hbt!] 
\centering
    \includegraphics[width=0.9\linewidth]{images/480hzcurve.png}    
    \caption{Complex graph of DKD with one component where $f^* = 480$Hz aligns with the sample rate.}
    \label{fig:complexcoherent}
\end{figure}

\newpage
To further visualize the process of DKD, we can add an additional frequency with a different phase at a nearby offset. Importantly, the DKD algorithm's ability to distinguish between frequencies holds even for two that occupy the same frequency bin. The reason for this comes from the choice of initial conditions when performing our search. Given that we fix the initial phase to that of the bin, this phase will be the weighted average of the two components within that bin. As such, it does not matter which frequency DKD finds first, as it will be subtracted out and there will only be one frequency in that bin during the second iteration. 
\begin{figure}[hbt!] 
\centering
    \includegraphics[width=0.9\linewidth]{images/MultipleSines.png}    
    \caption{Complex graph of DKD with two components}
    \label{fig:multiple_sines}
\end{figure}
\newpage
\subsubsection{Spectrally Rich Signals}
For more complicated signals, DKD requires substantially more iterations to fully capture the underlying frequencies. Although this is expected behavior, small mismatch errors can occur and the algorithm must run additional times to compensate. Nonetheless, after the algorithm has terminated many of these components will have small amplitudes and can be either filtered out or ignored depending on what is required. Furthermore, because DKD starts with frequencies of higher magnitudes, we know that for most real-world signals, lower frequencies tend to have much higher magnitudes than higher frequencies. This is especially true for audio signals with harmonics~\cite{de_Boer}. As a consequence, DKD tends to work from low frequencies to high for real-world audio as the number of iterations increases. \\
\begin{figure}[hbt!] 
\centering
    \includegraphics[width=0.9\linewidth]{images/spectral_waterfall.png}    
    \caption{Higher frequencies are underrepresented at low component counts}
    \label{fig:waterfall}
\end{figure}
\newline
For the figures in this section, we use snippets of ``Kidding Myself,'' a song from a folk-country band Sure Thing composed entirely of Brown students (including the author of this paper). We feel that this audio is a good choice for the purposes of visualization, as the song is composed of quality recordings of multiple instruments with rich timbres and minimal distortion. Qualitatively, as the low-component decompositions are still 48kHz, the reconstructions do not sound under-sampled; rather, they sound more like low-pass filtered versions of the original audio. It is also worth noting that there is some high frequency noise at medium component counts (2k-10k components) for complicated audio, this is likely because at mid to high frequencies, the effect of small (phase) errors in the kernel deconvolutions are perceptually very significant. Furthermore, as DKD essentially approximates the set of continuous frequencies that make up a real world signal, it will almost always result more frequencies than there are peaks in the DFT. Depending on the complexity of the signal, the algorithm can produce a set of frequencies that appear more ``noisy'' than the underlying DFT would suggest, especially for low sample rates or short signals.

It is worth noting that this behavior is desired, but the sorting step can be made more sophisticated to avoid the implicit LP filtering depending on the application. We experimented with this but added emphasis on mid frequencies without knowing the kernel effects of lower frequencies with even a small amplitude and produce the incorrect phase result, which is discussed further in Section \ref{sec:dkdni}. Ultimately, addressing this issue for the iterative algorithm appears to be extremely challenging and goes beyond the scope of the paper. 
\begin{figure}[hbt!]
    \includegraphics[scale=0.275]{images/stft_comparison.png}
    \includegraphics[scale=0.275]{images/spectral_difference.png}
    \caption{(left) STFT spectrogram of audio at various component counts, original at top; (right) difference in spectrograms at different component levels}
\end{figure}

\clearpage
To quantify the algorithm's overall performance both in terms of runtime and estimation error of the reconstructed signal $\bf{\hat{S}}$ and the original signal $\bf{S}$, we use several key metrics 
\begin{itemize}
    \item $L_2$ Convergence = $\displaystyle\sum_{i=0}^{N-1} \left|\left|DFT(\mathbf{S})_i - DFT(\mathbf{\hat{S}})_i \right|\right|_2$ \\
    \item Mean Square Error (MSE) = $\displaystyle\frac{1}{N}\sum_{i=0}^{N-1} \left|\left|DFT(\mathbf{S})_i - DFT(\mathbf{\hat{S}})_i\right|\right|_2^2$ \\
    \item Root Mean Square Error (RMSE) = $\displaystyle\sqrt{\frac{\sum_{i=0}^{N-1} \left|\left|DFT(\mathbf{S})_i - DFT(\mathbf{\hat{S}})_i\right|\right|_2^2}{N}}$ \\
    \item Cosine Similarity ($S_C$) = $\displaystyle\frac{\sum_{i=0}^{N-1}\Re\left( DFT(\mathbf{S})_i \cdot
 \overline{DFT(\mathbf{\hat{S}})_i}\right)}{\sum_{i=0}^{N-1}\left|\left|DFT(\mathbf{S})_i \right|\right|\cdot \sum_{i=0}^{N-1}\left|\left|DFT(\mathbf{\hat{S}})_i \right|\right|}$\\
    \item Log-Spectral Distance (LSD) = $\displaystyle\frac{1}{N}\sum_{i=0}^{N-1}\log_{10}DFT(\mathbf{S})_i - \log_{10}DFT(\mathbf{\hat{S}})_i$ 
    
\end{itemize}

We use the $L_2$ norm of the difference between the signals to determine how fast the reconstructed spectrum converges to that of the original. RMSE is a similarly good metric that is commonly used in speech processing. Cosine similarity is used to measure how similar the relative spectra are, with emphasis placed on the frequency and phase rather than the amplitudes. Finally, log-spectral distance is used as a log power scaled perceptual metric that mimics human sensitivity to relative magnitudes at different frequencies~\cite{lsd}.

Note that in all tests corresponding to signal length, the added time corresponded to new information (additional audio segments). We did this in part to differentiate from sample rate tests, but also to show how the algorithm performs in a more realistic scenario.
\newpage
\subsection{Runtime Analysis}

\subsubsection{Runtime Comparison Across Threads and Signal Length}

Although the FFT is famously an $\mathcal{O}(n\log n)$ operation, the constant term in front of the linear time main loop dwarfs that of the initial FFT. Thus, running the algorithm at any number of threads yields an approximately linear runtime with an $R^2$ value of $0.999$ for all five thread configurations. 
\begin{figure}[hbt!]
\centering
    \includegraphics[scale=0.5]{images/wall_time_vs_duration.png}
    \caption{5-Run average wall time for DKD with 1, 2, 4, 8, and 12 threads.}
\end{figure} 

\subsection{Convergence}
\subsubsection{Convergence Rate vs Signal Length}
By far the largest contributor to both the speed and information captured per iteration is the initial length ($N$) of the signal. This is primarily because longer signals increase the number of bins the algorithm must consider when updating the residual spectrum. Looking at the $L_2$ norm of the total energy contained in the residual relative to the original, we can see that the algorithm is not only convergent, but that the actual value approaches zero fairly quickly for lower component counts. As such, nearly all further graphs will be log-scaled. Regardless, the figure below shows how convergence is comparatively much slower for longer signals. 
\begin{figure}[hbt!]
\centering
    \includegraphics[scale=0.4]{images/convergence_rate.png}
    \caption{\% $L_2$ norm remaining in the residual versus the original spectrum}
\end{figure} 

We also wanted to evaluate the relative distance of the signals if we just directly stored the top $n$ largest FFT bins, which can be seen plotted as dashed lines. Although all metrics generally tend towards zero as component count increases, LSD remains better for DKD for nearly all component counts.
\begin{figure}[hbt!]
\centering
    \includegraphics[scale=0.65]{images/lsd_comparison.png}
    \caption{Log-spectral distance comparison across samples of varying lengths}
\end{figure} 

\begin{figure}[hbt!]
    \centering
    \includegraphics[width=0.8\linewidth]{images/accuracy_metrics_grid.png}
    \caption{$S_C$, LSD, and RMSE at multiple signal lengths}
    \label{fig:grid}
\end{figure}


\subsubsection{Convergence Rate vs Sample Frequency}
Initially, we believed that sample rate and signal length would exhibit similar returns to scale--especially on perceptual metrics. What we found however is that the algorithm is still fairly robust to under-sampled audio. However for 8kHz and 11kHz, we observed that DKD struggled a bit to fully capture all frequencies. This was likely because the lower sample rate but constant length force more high-amplitude frequencies to occupy the same bins, which is harder for DKD to work with. Furthermore, aliasing can also negatively impact DKD performance, as artifacts can disrupt the correlation calculation if the aliasing is severe enough. Specifically, we can see in Figures \ref{fig:sr_mse} and \ref{fig:sr_lsd} that at 15k-20k components, the 8kHz and 11kHz samples plateaued in both MSE and LSD.

\begin{figure}[hbt!]
    \centering
    \includegraphics[width=0.9\linewidth]{images/sample_rate_mse.png}
    \caption{Mean Squared Error at multiple sample rates}
    \label{fig:sr_mse}
\end{figure}
\begin{figure}[hbt!]
    \centering
    \includegraphics[width=0.9\linewidth]{images/sample_rate_lsd.png}
    \caption{Log-spectral distance at multiple sample rates}
    \label{fig:sr_lsd}
\end{figure}


Overall, the general trend we observed is that as sample rate increases, DKD can theoretically capture more components. This is partly because it is easier for the algorithm to pick out relevant true frequencies when bins are narrower, but also because there are simply more frequencies allowed with a higher Nyquist rate.



\chapter{Dirichlet Kernel Deconvolution: Applications}
\label{ch:5}
\section{Idealized Frequency Domain Resampling}
\label{sec:resample}

\subsection{Motivation and Background}

Resampling is the process of changing a signal's sampling rate while preserving its content as faithfully as possible. Traditional approaches to resampling operate in the time domain, typically using sinc-interpolation techniques or polyphase filters that can introduce artifacts such as aliasing, imaging, and loss of high-frequency information. Even with sophisticated windowing functions, time-domain resampling faces fundamental limitations when changing sampling rates, particularly when downsampling~\cite{RESAMPLE}.

The most common approach to high-quality resampling involves:
\begin{enumerate}
    \item Applying an LP anti-aliasing filter to remove frequencies above the new Nyquist limit
    \item Interpolating between existing samples to generate the new sample points
    \item Decimating (for downsampling) or inserting zeros and filtering (for upsampling)
\end{enumerate}

However, these approaches are all fundamentally constrained by the discrete time representation of the signal. Once high-frequency information is eliminated by anti-aliasing filters during downsampling, it cannot be recovered. Furthermore, interpolation between samples always represents an approximation to the true underlying continuous signal.

We propose a fundamentally different approach: performing resampling in the frequency domain using the Dirichlet Kernel Deconvolution (DKD) method introduced in Chapter \ref{ch:4}. This approach addresses several core limitations of time-domain resampling by working directly with the idealized frequency components rather than discretized time samples.

\subsection{Theoretical Foundation}

Our approach is based on the observation that a continuous signal can be represented as a sum of pure sinusoidal components, each with a specific frequency, amplitude, and phase. The DKD method allows us to decompose a discretely sampled signal into these components with sub-bin frequency precision. 

Let a sampled audio signal $x[n]$ with sample rate $f_{s1}$ be represented by a set of $M$ sinusoidal components:
\begin{equation}
\{(A_i, f_i, \varphi_i)\}_{i=1}^M
\end{equation}
where $A_i$, $f_i$, and $\varphi_i$ are the amplitude, frequency (in Hz), and phase of each component, respectively.

To resample this signal to a new sample rate $f_{s2}$, we need to retain the physical frequencies (in Hz) while recomputing their bin positions in the new sampling context. For a signal of length $N_1$ at the original sample rate, the relationship between bin number $b_1$ and frequency in Hz is:
\begin{equation}
f_{Hz} = b_1 \cdot \frac{f_{s1}}{N_1}
\end{equation}
For the new sample rate $f_{s2}$, we must determine the new bin position $b_2$ that corresponds to the same physical frequency:
\begin{equation}
b_2 = f_{Hz} \cdot \frac{N_2}{f_{s2}} = b_1 \cdot \frac{f_{s1}}{N_1} \cdot \frac{N_2}{f_{s2}}
\end{equation}
where $N_2$ is the length of the resampled signal. To maintain the original signal duration, $N_2$ must satisfy:
\begin{equation}
\frac{N_1}{f_{s1}} = \frac{N_2}{f_{s2}} \implies N_2 = N_1 \cdot \frac{f_{s2}}{f_{s1}}
\end{equation}
The key insight is that when resampling in the frequency domain using DKD, we must:
\begin{enumerate}
    \item Preserve the physical frequencies (in Hz) of all components
    \item Recalculate their bin positions based on the new sample rate
    \item Filter out components whose frequencies exceed the new Nyquist frequency
    \item Reconstruct the signal using these remapped components
\end{enumerate}

Crucially, unlike time-domain resampling where information above the new Nyquist frequency is irrecoverably lost, our approach perfectly preserves the parameters of all components within the new bandwidth limit. This allows for idealized resampling with minimal artifacts and maximum fidelity.

\subsection{Implementation}

Our implementation consists of the following steps:

\subsubsection{Decomposition Phase}
First, we apply the DKD algorithm to the original signal to obtain its component representation:

\begin{lstlisting}[language=C++, 
                  caption=Component Extraction for Resampling, 
                  basicstyle=\footnotesize\ttfamily,
                  breaklines=true]
// Extract components from original signal
std::vector<Component> components = deconvolveParallelDirichlet(fft, sample_rate, num_threads, maxcomp);
\end{lstlisting}

\subsubsection{Frequency Conversion}
Next, we convert each component's frequency from the original sampling context to the new one, filtering out components above the new Nyquist limit:

\begin{lstlisting}[language=C++, 
                  caption=Frequency Conversion for Resampling, 
                  basicstyle=\footnotesize\ttfamily,
                  breaklines=true]
std::vector<Component> valid_components;
const double new_nyquist = static_cast<double>(target_sr) / 2.0;

for (const auto &[true_frequency, true_phase, amplitude]: components) {
    // Convert from bin to physical frequency (Hz)
    const double freq_hz = true_frequency * static_cast<double>(original_sr) / static_cast<double>(N);
    
    // Only keep components below the new Nyquist frequency
    if (freq_hz < new_nyquist) {
        // Convert from Hz back to bin in the new sample rate context
        const double new_freq_bin = freq_hz * static_cast<double>(new_N) / static_cast<double>(target_sr);
        valid_components.push_back({new_freq_bin, true_phase, amplitude});
    }
}
\end{lstlisting}

\subsubsection{Reconstruction}
Finally, we reconstruct the signal at the new sample rate using the valid components:

\begin{lstlisting}[language=C++, 
                  caption=Signal Reconstruction at New Sample Rate, 
                  basicstyle=\footnotesize\ttfamily,
                  breaklines=true]
std::vector<double> resampled = decompressComponents(valid_components, new_N, DC, max_magnitude);
// Or using parallel processing:
// resampled = decompressComponentsParallel(valid_components, new_N, DC, max_magnitude, num_threads);
\end{lstlisting}

\subsection{Advantages and Limitations}

\subsubsection{Advantages}

\begin{enumerate}
    \item \textbf{Preservation of Frequency Content}: Unlike traditional resampling methods that may smear or distort frequency content, our approach maintains the exact frequencies, amplitudes, and phases of all components within the new Nyquist limit.
    
    \item \textbf{Minimal Aliasing}: By explicitly filtering components above the new Nyquist frequency, our method achieves optimal anti-aliasing without the need for filter design or windowing functions.
    
    \item \textbf{Phase Coherence}: Phase relationships between frequency components are preserved, maintaining the temporal characteristics and transients of the original signal.
    
    \item \textbf{Arbitrary Resampling Ratios}: The method supports any resampling ratio without degradation in quality, unlike traditional approaches which perform optimally only for specific ratios (e.g., powers of 2).
    
    \item \textbf{Perfect Reconstruction}: For signals that can be perfectly represented by their DKD components, our resampling method provides theoretically perfect reconstruction within the bandwidth constraints.
\end{enumerate}

\subsubsection{Limitations}

\begin{enumerate}
    \item \textbf{Computational Complexity}: The DKD-based resampling requires far more computation than traditional methods, particularly for long signals with complex spectral characteristics.
    
    \item \textbf{Memory Requirements}: Storing and manipulating the component representation requires more memory than direct time-domain resampling.
    
    \item \textbf{Approximation Quality}: For signals that are not perfectly represented by a finite number of DKD components, there will be some approximation error in the resampled output. To do this in a perceptually indistinguishable way from traditional methods, DKD requires approximately $30\text{k}-50\text{k}\times N$ components for a complex, multi-source real-world $N$ second audio segment. 
    
    \item \textbf{Loss of Information Above Nyquist}: Like all resampling methods, frequencies above the new Nyquist limit cannot be preserved. However, unlike traditional methods, our approach provides a clean cutoff without introducing artifacts. Furthermore, we can keep track of the removed components if needed.
\end{enumerate}

\subsection{Experimental Results}

To evaluate our frequency domain resampling method, we conducted experiments comparing it with traditional resampling techniques across various types of audio content. 

The results demonstrate that our method preserves the amplitude and phase characteristics of frequency components more accurately than traditional approaches, particularly for signals with complex harmonic content such as musical instruments and speech. Listening tests confirm that the DKD resampled audio maintains better transient response and timbral quality, especially when significant downsampling is performed. Although it performs slightly worse on RMSE and SI-SDR, it leads the four methods on LSD. 

Table~\ref{tab:resample_metrics} shows quantitative metrics comparing our method with traditional resampling approaches.\footnote{Windowed-Sinc corresponds to librosa.resample \text{res\textunderscore type='sinc\textunderscore best'}; Polyphase Filter to  res\textunderscore type='polyphase'; and Kaiser to res\textunderscore type=\text{'kaiser\textunderscore best'}} Each resampling method is downsampled then upsampled back to 48kHz and compared with the original:

\begin{table}[h]
\centering
\begin{tabular}{|l|c|c|c|c|}
\hline
\textbf{Metric} & \textbf{DKD Resampling} & \textbf{Windowed-Sinc} & \textbf{Polyphase Filter} & \textbf{Kaiser}\\
\hline
RMSE & 0.00645 & 0.00612 & 0.00610 & 0.00625\\
\hline
SI-SDR & 27.281 & 27.742 & 27.779 & 27.552 \\
\hline
LSD (dB) & 1.661 & 4.133 & 2.092 & 3.664 \\
\hline
\end{tabular}
\caption{Comparison of resampling methods using Root Mean Square Error (RMSE) (lower is better), Scale-Invariant Signal to Distortion Ratio (SI-SDR) (higher is better), and log-spectral distance (LSD) (lower is better) metrics when resampling from 48kHz to 16kHz then resampling back to 48kHz.}
\label{tab:resample_metrics}
\end{table}

\begin{table}[h]
\centering
\begin{tabular}{|l|c|c|c|c|}
\hline
\textbf{Metric} & \textbf{DKD Resampling} & \textbf{Windowed-Sinc} & \textbf{Polyphase Filter} & \textbf{Kaiser}\\
\hline
RMSE & 0.00071 & 0.00043 & 0.00044 & 0.00067 \\
\hline
SI-SDR & 46.474 & 50.832 & 50.692 & 47.002 \\
\hline
LSD (dB) & 0.916 & 2.802 & 0.993 & 2.439 \\
\hline
\end{tabular}
\caption{Comparison of resampling methods using root mean square error (RMSE), scale-invariant signal to distortion ratio (SI-SDR), and log-spectral distance (LSD) metrics when resampling from 48kHz to 32kHz then resampling back to 48kHz.}
\label{tab:resample_metrics_32}
\end{table}
\newpage
\subsection{Applications}

Frequency domain resampling using DKD is particularly valuable in several application areas:

\begin{enumerate}
    \item \textbf{Audio Production}: High-quality sample rate conversion for professional audio with minimal artifacts.
    
    \item \textbf{Archival and Restoration}: Idealized downsampling of high-resolution archival recordings for distribution formats.
    
    \item \textbf{Virtual Instruments}: Perfect pitch-shifting and time-stretching without the spectral artifacts commonly associated with traditional methods.
    
    \item \textbf{Scientific Signal Processing}: Precise resampling of measurement data where preserving frequency characteristics is critical, such as in seismology, radio astronomy, and medical signal processing.
\end{enumerate}

The DKD-based frequency domain resampling represents a novel approach to resampling. By decomposing signals into their constituent frequency components with sub-bin precision and reconstructing them at the new sample rate, we achieve resampling that preserves the spectral characteristics of the original signal with unprecedented perceptual fidelity. 

Although more work is needed on computational optimization for DKD to be used in real-time or in place of traditional resampling methods, the technique presented here not only builds upon the theoretical foundation for existing methods but also demonstrates the broader utility of the DKD algorithm beyond spectral analysis, showing its potential as a tool in digital signal processing.

\section{Non-Iterative DKD}
\label{sec:dkdni}
\subsection{Motivation}
Recall that the purpose of deconvolving the Dirichlet kernel is twofold: eliminating spectral leakage and decomposing the DFT into the set of continuous components that make up the signal. However, depending on that signal, these goals might be at-odds. Given that in most real-world signals there are many more component frequencies than there are DFT bins, DKD could result an ostensibly messier spectrum than simply looking at the original rectangular-windowed DFT. 

The fundamental challenge to this type of deconvolution is ensuring correct phases, which is why DKD must be run iteratively for proper reconstruction. Let us consider running DKD across all peaks in the DFT concurrently. Although this would produce a relatively accurate approximation to the true frequencies for the subset of components that most strongly contribute to each peak bin, the phases returned by the algorithm might not correspond to the true phases of the isolated components. This is because the returned ``optimal'' phases would still be affected by the spectral leakage from all other kernels. In other words, the optimal phases would be ``correct'' in the sense that they correspond to the frequency-adjusted peak given the current bin's phase, but incorrect in approximating the true underlying phase of the actual continuous sinusoid. Moreover, as there are often far more real-world components than there are DFT peaks, the returned set of frequencies would not fully encapsulate the true set of components.

The largest practical consequence of this is that DKD must be run iteratively from highest to lowest magnitudes in order to ensure that the large components, which have far-reaching kernel effects on less important frequencies (and/or higher frequencies with intrinsically lower magnitudes) dwarfing the relative magnitudes of the smaller kernels, are not considered when the algorithm computes the phases of the latter set of frequencies. Importantly, the first set of frequencies is not immune to this same type phase issue, which is actually where much of the approximation error comes from (and why noise is audible at low-medium component counts). As DKD iterations increase, the algorithm inserts minor ``virtual'' frequencies near larger peaks to account for small phase mismatches, but this behavior is effective only when these differences are small. Hence, we must ensure that each component is considered only with kernel interference from less important frequencies.\footnote{While it is convenient that DKD is robust to these types of errors, there is likely a way to consolidate the redundant components by adjusting the phases of nearby larger ones. This goes beyond the scope of the paper.}

Although DKD's iterative approach allows for robustness to minor phase issues, it also imposes serious limitations on computational efficiency. That said, there are cases in which reconstruction and phase-faithful decomposition are less important than a clean visualization or a good estimate of the most important continuous frequencies. In these cases, we propose a modification to the original DKD algorithm in which we only consider the DFT peaks and run the algorithm non-iteratively. This also offers a tremendous advantage to computational efficiency, as we can parallelize the algorithm to evenly distribute each peak across all cores. Furthermore, since we do not need to update the residual if DKD is performed in one pass, we save on the majority of the actual computation time and completely avoid compounding effects from waiting on threads.

\subsection{Algorithm}

The non-iterative approach to Dirichlet Kernel Deconvolution significantly modifies the control flow of the original algorithm. While the mathematical foundation remains the same, the algorithm no longer performs sequential extraction and subtraction of components. Instead, all significant peaks in the spectrum are identified in a single pass, and their corresponding frequency, phase, and amplitude parameters are optimized in parallel.

\subsubsection{Algorithm Overview}

The non-iterative DKD algorithm follows these steps:

\begin{enumerate}
    \item Identify all significant peaks in the magnitude spectrum (local maxima)
    \item Sort the peaks by magnitude in descending order
    \item Process each peak independently in parallel to find its optimal frequency, phase, and amplitude
    \item Collect the resulting components as the final decomposition
\end{enumerate}

This approach forgoes the iterative refinement of the residual spectrum, trading some reconstruction accuracy for significant computational efficiency. The key insight is recognizing that for many applications, a faithful identification of the principal spectral components is more important than the perfect isolation of every component's phase.

\subsubsection{Peak Identification}

Unlike the iterative algorithm, which processes one peak at a time, the non-iterative version first identifies all significant peaks in the spectrum:

\begin{lstlisting}[language=C++, 
                   caption=Peak Identification in Non-Iterative DKD, 
                   basicstyle=\footnotesize\ttfamily,
                   breaklines=true,
                   breakatwhitespace=true]
// Step 1: Find all significant peaks across the spectrum
std::vector<PeakInfo> all_peaks;
for (size_t i = 1; i < nyquist_bin; i++) {
    double magnitude = std::abs(fft[i]);

    // Simple peak detection - local maximum with meaningful magnitude
    if (i > 1 && i < nyquist_bin &&
        magnitude > std::abs(fft[i-1]) &&
        magnitude > std::abs(fft[i+1])) {

        PeakInfo peak{};
        peak.bin = i;
        peak.magnitude = magnitude;
        peak.frequency = i; // Initial guess
        peak.phase = std::arg(fft[i]);
        all_peaks.push_back(peak);
    }
}

// Sort peaks by magnitude (descending)
std::sort(all_peaks.begin(), all_peaks.end(),
    [](const PeakInfo& a, const PeakInfo& b) {
        return a.magnitude > b.magnitude;
    }
);

// Limit total number of peaks to process
if (all_peaks.size() > MAX_COMPONENTS) {
    all_peaks.resize(MAX_COMPONENTS);
}
\end{lstlisting}

This approach ensures that the algorithm prioritizes the most energetically significant components while still maintaining a comprehensive view of the spectrum.

\subsubsection{Parallel Component Optimization}

The most substantial improvement in the non-iterative algorithm comes from parallelizing the component parameter optimization. Each peak is processed independently by a dedicated thread:

\begin{lstlisting}[language=C++, 
                   caption=Parallel Component Optimization, 
                   basicstyle=\footnotesize\ttfamily,
                   breaklines=true,
                   breakatwhitespace=true]
// Step 2: Process these peaks in parallel
std::vector<Component> all_components;
std::mutex components_mutex; // For safely appending components

// Create thread pool
std::vector<std::thread> threads;

// Calculate peaks per thread
size_t peaks_per_thread = (all_peaks.size() + num_threads - 1) / num_threads;

for (size_t t = 0; t < num_threads; t++) {
    size_t start_idx = t * peaks_per_thread;
    size_t end_idx = std::min((t + 1) * peaks_per_thread, all_peaks.size());

    if (start_idx >= all_peaks.size()) continue;

    threads.emplace_back([&, start_idx, end_idx]() {
        std::vector<Component> thread_components;

        for (size_t i = start_idx; i < end_idx; i++) {
            const PeakInfo& peak = all_peaks[i];

            // Skip if peak frequency is too close to 0 or Nyquist
            if (peak.bin < 2 || peak.bin > nyquist_bin - 2) continue;

            // Process this peak
            try {
                // Find optimal frequency using golden section search
                double true_freq = findOptimalFrequency(fft, peak.bin, peak.bin);

                // Find optimal phase
                double true_phase = findOptimalPhase(fft, peak.bin, true_freq);

                // Determine true amplitude
                double amplitude = DirichletKernel::getAmplitudeAtBin(
                    true_freq, std::abs(fft[peak.bin]), N, peak.bin);

                thread_components.push_back({true_freq, true_phase, amplitude});
            }
            catch (const std::exception& e) {
                std::cerr << ``Error processing peak at bin `` << peak.bin
                          << ``: `` << e.what() << std::endl;
            }
        }

        // Add this thread's components to the global list
        std::lock_guard<std::mutex> lock(components_mutex);
        all_components.insert(all_components.end(),
                             thread_components.begin(),
                             thread_components.end());
    });
}

// Wait for all threads to complete
for (auto& thread : threads) {
    if (thread.joinable()) thread.join();
}
\end{lstlisting}

This parallel approach divides the peaks among available threads and processes them concurrently, dramatically reducing computation time.

\subsubsection{Parameter Optimization Functions}

The functions used to find optimal frequency, phase, and amplitude remain the same as in the iterative approach, ensuring mathematical consistency. The optimization process still uses golden section search to find the precise frequency and phase.

The key distinction is that these functions now operate on the original spectrum rather than on an iteratively updated residual, which means they find parameters that best fit the observed data without accounting for the influence of other components.

\subsubsection{Theoretical Impact on Accuracy}

The mathematical implication of the non-iterative approach deserves special attention. Let the original spectrum $\mathbf{F}$ be composed of $M$ sinusoidal components:

\begin{equation}
\mathbf{F} = \sum_{i=1}^{M} A_i D(f_i, \varphi_i)
\end{equation}

where $D(f_i, \varphi_i)$ represents the Dirichlet kernel response for a component with frequency $f_i$ and phase $\varphi_i$.

In the iterative approach, when extracting the $j$-th component, we compute parameters that best fit the residual:

\begin{equation}
\mathbf{R}_j = \mathbf{F} - \sum_{i=1}^{j-1} A_i^* D(f_i^*, \varphi_i^*)
\end{equation}

By contrast, in the non-iterative approach, each component is optimized against the complete original spectrum:

\begin{equation}
(A_j^*, f_j^*, \varphi_j^*) = \argmax_{A, f, \varphi} \text{Correlation}(D(f, \varphi), \mathbf{F})
\end{equation}

This distinction means that components extracted non-iteratively will have parameters that account for the spectral leakage from all other components, potentially leading to phase estimates that differ significantly from those that would be obtained through iterative extraction.

\subsubsection{Computational Complexity Analysis}

The computational complexity of the non-iterative algorithm is:

\begin{equation}
O\left(\frac{P \cdot \ell \cdot \log(K)}{T}\right)
\end{equation}

where:
\begin{itemize}
    \item $P$ is the number of peaks processed
    \item $\ell$ is the correlation window size
    \item $K$ is the precision of the frequency/phase search
    \item $T$ is the number of threads
\end{itemize}

This represents a significant improvement over the iterative approach's complexity of $O(M \cdot N \cdot \ell \cdot \log(K))$, particularly when $P \ll M \cdot N$ and multiple threads are available.

\subsubsection{Algorithm Limitations}

The principal limitation of the non-iterative approach is its inability to accurately reconstruct the original signal due to phase interaction effects. When multiple frequency components contribute significantly to the same spectral region, their phase relationships become entangled. Without the iterative extraction and residual updating process, these relationships cannot be disentangled, leading to phase estimates that optimize correlation with the observed spectrum but do not reflect the true underlying component phases.

This limitation makes the non-iterative algorithm less suitable for perfect reconstruction applications but maintains its utility for applications where the identification of principal frequency components is the primary goal, such as spectral analysis, feature extraction, and visualization.

\subsection{Results}
Qualitatively, the reconstructed audio from non iterative DKD sounds surprisingly coherent (albeit somewhat ``ethereal''). In fact, the high frequency information in Figure \ref{fig:stft_ni} is not random harsh noise; rather, it sounds more like a lush reverberation effect. Most of the underlying information is present in the reconstruction, but there are clearly some frequencies missing in more complex audio--almost certainly cases where there are multiple true frequencies to a bin. What we found most surprising is that the difference between the non iterative reconstruction and the original audio is smaller as the input gets larger. This is likely a result of added bin resolution, which means that the DFT peaks more closely correspond to the real number of true frequencies, and because the added space between peaks reduces the spectral leakage across bins. \\

The runtime difference for decomposition is staggering. There is an 7200\%-52500\% overall speedup when both the iterative and non-iterative are run with \texttt{max\_components} equal to the initial number of DFT peaks using 12 threads. \\

This difference also means that non-iterative DKD can be feasibly run in real-time, which opens the doors for potential applications that do not require a faithful reconstruction. It is important to note that the table above is only the time it takes to decompose the DFT into a set of true frequency, phase, and amplitude components and does not include reconstruction time, which is generally comparable to the iterative decomposition time.\\

\begin{table}[h]
\centering
\begin{tabular}{|l|c|c|c|c|}
\hline
\textbf{Length (s)} & \textbf{Peak ct.} & \textbf{Wall Time (s) - Non-Iterative} & \textbf{Wall Time (s) - Iterative} \\
\hline
2 & 16369 &  0.110 & 7.941 \\
\hline
4 & 28592 & 0.183 & 23.869 \\
\hline
8 & 61548 & 0.364 & 96.319 \\
\hline
16 & 131226 & 0.756 & 397.050 \\
\hline
\end{tabular}
\end{table}

\begin{figure}[hbt!]
    \centering
    \includegraphics[width=1.0\linewidth]{images/stft_2s_corrected.png}
    \caption{2 second STFT spectrogram comparison}
    \label{fig:stft_ni}
\end{figure}
\clearpage
\begin{figure}[hbt!]
    \centering
    \includegraphics[width=1.0\linewidth]{images/stft_4s.png}
    \caption{4 second STFT spectrogram comparison}
    \includegraphics[width=1.0\linewidth]{images/stft_ni_8s_corrected.png}
    \caption{8 second STFT spectrogram comparison}
\end{figure}

\clearpage
\section{Application: Cleaned Short-Time Fourier Transform}
\subsubsection{Motivation}
\label{subsec:cstft_intro}
The largest drawback to using the Fourier transform as an analysis tool is that one must completely trade resolution in the time domain for resolution in the frequency domain~\cite{smith2011spectral}. In other words, although the DFT reveals much about the underlying periodic frequencies in the data, it gives no information as to when those frequencies might appear. To address this, we can use the Short-Time Fourier Transform (STFT), which works by dividing the input signal into $\texttt{n\_fft}$ sized segments with some constant overlap using an offset equal to some other (smaller) number of samples--called the hop size and denoted by $\texttt{hop\_length}$. Then, we would traditionally apply some window function\footnote{Common windows are Hann, Hamming, etc.} to smooth out the ends of the segments and avoid artifacts. Finally, we take the DFT of each of these segments and stitch them together to get a ``best of both worlds'' analysis on the signal's spectrum over time. 

Although time-frequency analysis is essential in signal processing, traditional STFT approaches suffer from some fundamental limitations. Specifically, there are constraints on the choice of $\texttt{hop\_length}$ and $\texttt{n\_fft}$~\cite{smith2011spectral} which prevent the STFT from resolving closely-spaced frequencies with high resolution--especially when a signal is rich. As the length of each segment must be incredibly small, the resulting STFT bins will be fairly wide. Thus, each individual slice will have far worse resolution than a DFT of the full signal. Consequently, STFT spectrograms can often look messy due to a high degree of spectral leakage and frequency overlap caused by the wider bins accumulating too much spectral information, which can ultimately lead to certain frequencies blending with others. 

Having developed an effective solution to spectral leakage through Dirichlet Kernel Deconvolution, we can apply this technique frame-by-frame to create a ``Cleaned STFT.'' Since spectral leakage affects each STFT frame individually, applying DKD to each frame allows us to extract the true underlying frequencies with sub-bin precision, while still maintaining the temporal information provided by the STFT structure.

The key insight is that although traditional STFT is constrained by the discrete frequency bins determined by the FFT size, our DKD approach allows each frame to be decomposed into components with continuous frequency values. This means we can achieve frequency precision far beyond what the bin width would normally allow, while still preserving the time-domain information that makes STFT valuable.

Furthermore, by eliminating spectral leakage in each frame, we significantly reduce the ``smearing'' effect that obscures closely spaced frequency components in traditional spectrograms. This is particularly valuable for analyzing complex audio signals with rich harmonic content, where the true frequency components can be masked by leakage from stronger components nearby.

In this section, we present a complete implementation of this Cleaned STFT approach, demonstrate its advantages compared to traditional STFT through visual examples, and discuss its applications and limitations. We show that by combining the frame-wise processing of STFT with the spectral precision of DKD, we can achieve a time-frequency representation that more accurately reflects the true content of complex signals.

Moreover, since STFT spectrograms do not rely on phase information,\footnote{Although STFT spectrograms only show magnitude, phase information is often required for STFT tasks woutside of visualization. However, we can derive STFT phases simply from the magnitudes using the Griffin-Lim Algorithm~\cite{Hasegawa_2023}} visualizing the Cleaned STFT is a perfect application for Non-Iterative DKD. This allows us to take advantage of the dramatic improvements to computational efficiency without worrying about phasing issues.
\subsubsection{Theoretical Foundations}
The traditional STFT maps a one-dimensional signal $x(t)$ to a two-dimensional time-frequency representation $S(\tau, \omega)$ by computing:

\begin{equation}
S(\tau, \omega) = \int x(t)w(t - \tau)e^{-j\omega t} dt
\end{equation}

Where $w(t)$ is a window function centered at time $\tau$. In discrete implementation, this becomes a series of FFTs applied to windowed segments of the signal. Each segment is processed independently, producing a column in the resulting spectrogram.

The key limitation of this approach is that each segment's FFT produces discrete frequency bins with a resolution directly determined by the segment length. Specifically, the frequency resolution $\Delta f = f_s/N$, where $f_s$ is the sample rate and $N$ is the segment length. This quantization forces all frequency components to be represented at the nearest bin, causing information loss and contributing to the ``blurry'' appearance of traditional spectrograms.

Our Cleaned STFT approach replaces the FFT step with Dirichlet Kernel Deconvolution, maintaining the same segmentation and windowing structure but extracting continuous frequency components from each frame. Mathematically, for each frame $i$, we transform the windowed segment $x_i(t)$ into a set of components $\{(A_j, f_j, \phi_j)\}$ where:

\begin{itemize}
    \item $A_j$ represents the amplitude of component $j$
    \item $f_j$ represents the precise frequency of component $j$ (not constrained to FFT bins)
    \item $\phi_j$ represents the phase of component $j$
\end{itemize}

The resulting time-frequency representation can be expressed as:
\begin{equation}
C(\tau, f) = \sum_j A_j(\tau)\delta(f - f_j(\tau))e^{j\phi_j(\tau)}
\end{equation}

Where $\delta$ is the Dirac delta function. This representation preserves the temporal structure of the STFT while providing continuous frequency resolution within each frame.

\subsubsection{Implementation}
Our implementation of the Cleaned STFT follows these steps:

\begin{enumerate}
    \item \textbf{Frame Generation}: The input signal is divided into overlapping frames of length \texttt{n\_fft} with a hop size of \texttt{hop\_length} samples. Each frame is multiplied by a window function\footnote{We use the Hann window in our experiments. As it turns out, the CSTFT performs better when a window function other than rectangular is used. It is unclear why this is the case.} to reduce edge effects.
    
    \item \textbf{Component Extraction}: Rather than computing an FFT for each frame, we apply the non-iterative DKD algorithm to extract frequency components with sub-bin precision. This process identifies the true frequencies present in each frame, along with their amplitudes and phases.
    
    \item \textbf{Time-Frequency Mapping}: The extracted components from each frame are mapped to their corresponding time positions, creating a continuous time-frequency representation.
    
    \item \textbf{Visualization}: For visualization purposes, we bin the components to create a traditional-looking spectrogram (sacrificing precision at lower frequencies but better for facilitating comparison).
    
\end{enumerate}

The key parameters affecting the performance of our Cleaned STFT are:

\begin{itemize}
    \item \textbf{Frame Size (\texttt{n\_fft})}: Determines the basic frequency resolution and time localization tradeoff
    \item \textbf{Hop Length}: Controls the time resolution of the resulting spectrogram
    \item \textbf{Maximum Components}: Limits the number of frequency components extracted per frame, balancing detail against computational efficiency
    \item \textbf{Window Function}: Affects the spectral characteristics of each frame before processing
\end{itemize}

For our experiments, we use \texttt{n\_fft=2048}, \texttt{hop\_length=512}, and \texttt{max\_components=20000} with a Hann window, which provides a good balance of frequency resolution and temporal detail for typical audio signals sampled at 48kHz. We choose a purposefully high value for \texttt{max\_components}, as we want to ensure it is greater than the identified number of peaks.

\subsubsection{Experimental Results}

Figure \ref{fig:stft_waves_col} and \ref{fig:stft_comparison_log} show comparisons between the traditional STFT and our Cleaned STFT approach applied to an 8-second audio sample. The traditional STFT (top) exhibits typical spectral leakage, frequency overlap, and overall limited component resolution, while the Cleaned STFT (bottom) provides sharper frequency localization and clearer harmonic structures. 
\clearpage
Several key features are specifically worth highlighting:
\begin{itemize}
 \item \textbf{Temporal Dynamics: } Finally, the cleaned STFT preserves the temporal dynamics of the input signal without the presence of large leakage patterns after each impulse. In Figure \ref{fig:stft_comparison_log}, we see that note impulse responses in the traditional STFT are clear because they are spread across many different frequencies. Although this can be helpful visually, it can mask other frequencies and obfuscate more complex spectral dynamics associated with the ADSR envelope of a note. Cleaned STFT solves this problem by removing spectral leakage while still displaying clear differences in amplitude and harmonic content at new note impulses. This allows for both a better visualization of how frequencies near the main harmonics of a note change over time and a clearer picture of the overall temporal structure shown in the spectrogram.  
 
\item \textbf{Harmonic Structure: } The top/bottom left images in Figure \ref{fig:stft_waves_col} demonstrate the Cleaned STFT's improved ability to capture the subtle harmonic structure of a set of vocal overtones over time. While the traditional STFT shows the general shape of the vibrato and its effects on the overtone series, it does not clearly convey the specific frequencies at each time frame. However, the Cleaned STFT shows not only how the series changes with pinpoint accuracy, but it also captures a much clearer picture of the quieter frequencies hidden between the overtones that would otherwise get masked by spectral leakage. Furthermore, if we observe the difference in low-frequency resolution between the original and the Cleaned spectrograms, we see that while the traditional STFT displays a fairly muddled mix of frequencies, the Cleaned STFT maintains clarity and readability even with frequencies bunched together when displayed on a linear frequency axis. 
    
\item \textbf{Frequency Resolution: } The right side images in Figure \ref{fig:stft_waves_col} show how the Cleaned STFT not only offers dramatic improvements to frequency resolution, but enables tracking how partials evolve over time. The traditional STFT spectrogram has nowhere near enough precision to accomplish this task.\footnote{Audio Processing tools like SPEAR do allow for this currently, but they use more advanced methods than simply the STFT alone} Importantly, the approximated continuous frequencies displayed below are quantized back to fit in the traditional STFT's bin resolution for understandability, but the true frequencies obtained by DKD are even less jagged than the visualization would suggest. 
\end{itemize}

\begin{figure}

\centering
\begin{tikzpicture}
  
  \node[anchor=south west,inner sep=0] (image) at (0,0) 
       {\includegraphics[width=1.25\textwidth]{images/TEST_8s_comparison_log8192_hl128.png}};
  
\end{tikzpicture}
\caption{Cleaned STFT comparison with log-scaled frequency axis, \texttt{n\_fft} = 8192, \texttt{hop\_length} = 128. (Cleaned STFT quantized to match resolution of STFT)}
\label{fig:stft_comparison_log}
\end{figure}


\begin{figure}
\centering
\begin{tikzpicture}

  \node[anchor=south west,inner sep=0] (image) at (0,0) 
       {\includegraphics[width=1.2\textwidth]{images/TEST_8s_comparison_hz2048.png}};
  
  
  \begin{scope}[x={(image.south east)},y={(image.north west)}]
    
    \draw[red,thick,rounded corners] (0.06,0.0575) rectangle (0.15,0.1475);

    \draw[red,thick,rounded corners] (0.06,0.55) rectangle (0.15,0.65);

    \draw[green,thick,rounded corners] (0.69,0.175) rectangle (0.73,0.225);

    \draw[green,thick,rounded corners] (0.69,0.675) rectangle (0.73,0.725);
    
    
    \draw[red,thick,-stealth] (0.075,-0.03) -- (0.1,0.0575);
    \draw[red,thick,-stealth] (0.075, 1.03) -- (0.1,0.65);

    \draw[green,thick,-stealth] (0.75,-0.03) -- (0.71,0.175);
    \draw[green,thick,-stealth] (0.75, 1.03) -- (0.71,0.725);
    
    
    \node[anchor=north west] at (0.0,0.0) 
         {\includegraphics[width=0.55\textwidth]{images/waves_dkd.png}};
    \node[anchor=south west] at (0.0,1.0) 
         {\includegraphics[width=0.5\textwidth]{images/waves_stft.png}};
    \node[anchor=north west] at (0.5,0.0) 
         {\includegraphics[width=0.45\textwidth]{images/column_dkd.png}};
    \node[anchor=south west] at (0.5,1.0) 
         {\includegraphics[width=0.45\textwidth]{images/column_stft.png}};
    
  \end{scope}
\end{tikzpicture}
\caption{Comparison of Traditional STFT (top) and Cleaned STFT (bottom) with zoomed region highlighting the improved harmonic resolution achieved by the DKD approach. Hann window, \texttt{n\_fft}=2048, \texttt{hop\_length}=512}
\label{fig:stft_waves_col}
\end{figure}
\newpage
\subsubsection{Advantages}

\begin{enumerate}
    \item \textbf{Superior Frequency Resolution}: The Cleaned STFT resolves frequencies with continuous precision rather than aggregation into to the nearest FFT bin.
    
    \item \textbf{Reduced Spectral Leakage}: By applying DKD to each frame, we eliminate the spectral leakage that causes energy smearing and ``blurred'' visualizations in traditional spectrograms.
    
    \item \textbf{Improved Component Separation}: Closely spaced frequencies that would merge together in a traditional STFT can be clearly distinguished.
    
    \item \textbf{Preserved Temporal Structure}: Unlike some alternative high-resolution spectral techniques, our approach maintains the same temporal resolution as the traditional STFT.
    
    \item \textbf{Sparsity}: The component-based representation is naturally sparse, potentially enabling more efficient storage and processing compared to dense spectrograms.
\end{enumerate}

\subsubsection{Limitations}

\begin{enumerate}
    \item \textbf{Computational Cost}: The current implementation requires approximately 5-7x the processing time of traditional STFT, though this could be further optimized.\footnote{This figure is hard to calculate, as rendering the visualization through librosa and scheduling FFT jobs takes more time than actually running DKD. The 5-7x figure comes only from the time spent running DKD and librosa's STFT, although an optimized version would likely implement its own STFT in C++ rather than inter-operating between C++ and python using a bash script as we have done.}
    
    \item \textbf{Parameter Sensitivity}: The quality of results depends on appropriate selection of parameters, particularly \texttt{n\_fft} and the length of the file for both visualization purposes and for bin resolution.
    
    \item \textbf{Time-Frequency Tradeoff}: Like the traditional STFT, our approach still faces the fundamental tradeoff between time and frequency resolution determined by the frame size. Since single-pass NIDKD will never produce more frequencies then there are DFT peaks, the CSTFT will not be able to pull apart two frequencies in the same bin.
\end{enumerate}
\clearpage
\subsection{Further Applications of the CSTFT}

The Cleaned STFT offers several promising applications that would be worth exploring as potential future work in this area:

\subsubsection{Audio Analysis and Feature Extraction}
The improved frequency resolution enables more accurate analysis of musical content, including precise pitch tracking, harmonic analysis, and feature extraction for machine learning. By eliminating spectral leakage, the extracted features more accurately represent the true signal characteristics.

\subsubsection{Sound Transformation}
Time-stretching and pitch-shifting algorithms can benefit from the more accurate frequency representation, potentially reducing artifacts when modifying temporal or spectral characteristics.

\subsubsection{Source Separation/Identification}
The cleaner separation of frequency components may improve source separation algorithms, particularly for signals with closely spaced or overlapping harmonics. This Cleaned STFT approach could be merged with existing deep learning approaches on spectrograms to add additional precision and context that would not be available otherwise. This approach could potentially reduce model parameter counts substantially.

\subsection{Conclusion}

The Cleaned STFT represents a significant advancement in time-frequency analysis by combining the temporal structure of the traditional STFT with the frequency precision afforded by DKD. By extracting continuous frequency components from each frame, we achieve dramatically improved spectral clarity while maintaining temporal resolution.

Future work could additionally focus on optimizing the computational efficiency of the approach, developing specialized visualization tools for the continuous frequency representation, and applying the technique to specific audio analysis tasks such as music transcription or speech recognition.

\chapter{Frequency and Amplitude Modulation}
\label{ch:2}
A primary focus of this thesis is to analyze the properties and applications of frequency modulation in a more abstract mathematical context, as the interactions between pairs of FM operators exhibit interesting and useful behavior. First, we must define a ``simple'' FM operator array, which is defined by a carrier $c$, modulator $m$, modulation index $\beta$, and amplitude $A$.
\section{The Structure of an FM Synthesizer}
FM operators represent the fundamental building blocks of frequency modulation synthesis. As shown in the diagram, an FM synthesizer consists of four primary components:

\begin{enumerate}
    \item \textbf{Carrier Oscillator}: Produces a sine wave at frequency $c$, which serves as the base signal that will be heard.
    
    \item \textbf{Modulator Oscillator}: Generates a sine wave at frequency $m$, which is used to modulate the carrier's frequency.
    
    \item \textbf{Modulation Index ($\beta$)}: Controls the depth of modulation, determining how much the carrier frequency deviates from its center value. This parameter significantly affects the spectral richness of the output.
    
    \item \textbf{Amplitude ($A$)}: Scales the overall output signal volume.
\end{enumerate}

These components interact with a precise mathematical relationship. The modulator's output is scaled by the modulation index and then applied to the carrier's phase, resulting in a time-varying frequency. This interaction can be visualized as the modulator ``bending'' the carrier's frequency up and down according to its own oscillation.
\begin{figure}[hbt!]
    \centering
    \includegraphics[width=0.75\linewidth]{images/fm_waveform.png}
    \caption{FM waveform with increasing modulation index}
    \label{fig:fm_waveform}
\end{figure}

When analyzing how each parameter affects the spectral characteristics of an FM signal, we observe several key relationships~\cite{chowning1986fm}:

\begin{enumerate}
    \item \textbf{Carrier Frequency ($c$)}: Determines the perceived pitch of the sound and serves as the central frequency around which sidebands are distributed.
    
    \item \textbf{Modulator Frequency ($m$)}: Establishes the spacing between sidebands. When $m$ is harmonically related to $c$ (i.e., $m/c$ is a simple ratio), the resulting sound tends to be more harmonic and musical. Non-integer ratios produce more complex, potentially inharmonic spectra.~\cite{Chowning_1973}
    
    \item \textbf{Modulation Index ($\beta$)}: Controls the number and amplitude of sidebands. As $\beta$ increases:
    \begin{itemize}
        \item More energy is distributed to higher-order sidebands
        \item The carrier's amplitude decreases according to $J_0(\beta)$
        \item The spectral complexity increases dramatically
        \item At certain values of $\beta$ where $J_0(\beta) = 0$, the carrier disappears entirely
    \end{itemize}
    \item \textbf{Amplitude ($A$)}: Scales all spectral components uniformly without affecting their relative proportions.
\end{enumerate}

\begin{figure}[hbt!]
    \centering
    \includegraphics[width=0.7\linewidth]{images/fm_operator_diagram_emptybkg.png}
    \caption{Diagram of FM with two operators~\cite{Wikipedia_fmsvg}}
    \label{fig:fm_diagram}
\end{figure}
A critical feature of FM synthesis is how it handles spectral energy across the Nyquist frequency. When a sideband would mathematically extend beyond the Nyquist limit, it ``reflects'' back into the audible spectrum, potentially creating dissonant artifacts known as ``zero reflections'' or ``aliasing.'' These reflections follow the formula:
\begin{equation}
  f_{reflected} = 2f_{Nyquist} - f_{original}  
\end{equation}
This phenomenon must be carefully managed, particularly when using high modulation indices or high frequencies, to avoid undesired spectral components.\\
\begin{figure}[hbt!]
    \centering
    \includegraphics[width=0.9\linewidth]{images/FM_surface.png}
    \caption{Two operator FM spectrum as a function of modulation index~\cite{Chowning_1973}}
    \label{fig:fm_surface}
\end{figure}
Consider an FM operator array with amplitude $A$, carrier frequency $c$, modulator frequency $m$, and modulation index $\beta$. Let $FM(A, c, m, \beta)(t) = A \sin(ct + \beta \sin(mt))$ denote the signal generated by this arbitrary FM operator array.

Note that by its Fourier series, $\displaystyle FM(t) = A \sum_{n=-\infty}^\infty J_n(\beta) \sin((c + nm)t)$, where $J_n$ denotes the $n$-th Bessel function of the first kind.

\subsection{The Mathematical Connection Between FM and Bessel Functions}

The appearance of Bessel functions in the frequency spectrum of an FM signal is not arbitrary, but emerges naturally from the mathematical structure of frequency modulation. To derive this relationship, we begin with our standard FM expression:
\begin{equation}
  FM(A, c, m, \beta) = A \sin(ct + \beta \sin(mt))  
\end{equation}
To identify the frequency components present in this signal, we can rewrite it using Euler's identity:
\begin{equation}
  \sin(ct + \beta \sin(mt)) = \frac{e^{i(ct + \beta \sin(mt))} - e^{-i(ct + \beta \sin(mt))}}{2i}  
\end{equation}
Focusing on the exponential term, we can separate the carrier and modulator effects:
\begin{equation}
    e^{i(ct + \beta \sin(mt))} = e^{ict} \cdot e^{i\beta \sin(mt)}    
\end{equation}
Now, we can apply the Jacobi-Anger expansion, which gives us:
\begin{equation}
  e^{iz\sin\theta} = \sum_{n=-\infty}^{\infty} J_n(z)e^{in\theta}  
\end{equation}
where $J_n(z)$ is the $n$-th order Bessel function of the first kind. Applying this expansion to our FM signal with $z = \beta$ and $\theta = mt$:
\begin{equation}
  e^{i\beta \sin(mt)} = \sum_{n=-\infty}^{\infty} J_n(\beta)e^{inmt}  
\end{equation}
Combining with the carrier term:
\begin{equation}
  e^{i(ct + \beta \sin(mt))} = e^{ict} \cdot \sum_{n=-\infty}^{\infty} J_n(\beta)e^{inmt} = \sum_{n=-\infty}^{\infty} J_n(\beta)e^{i(c+nm)t}
\end{equation}
Substituting this back into our original expression and applying the same process to the negative exponential term:
\begin{equation}
  FM(A, c, m, \beta) = A\frac{1}{2i}\left(\sum_{n=-\infty}^{\infty} J_n(\beta)e^{i(c+nm)t} - \sum_{n=-\infty}^{\infty} J_n(\beta)e^{-i(c+nm)t}\right)  
\end{equation}
Using the relation $J_{-n}(\beta) = (-1)^n J_n(\beta)$ for integer values of $n$ and converting back to sine notation~\cite{kreh_2012}:\footnote{This is a rather involved set of steps and uses properties of Bessel functions coupled with trigonometric identities.}
\begin{equation}
  FM(A, c, m, \beta) = A\sum_{n=-\infty}^{\infty} J_n(\beta)\sin((c+nm)t)  
\end{equation}
This derivation reveals that an FM signal with modulation index $\beta$ decomposes into an infinite series of sinusoidal components. Each component has:
\begin{itemize}
    \item Frequency: $c + nm$ (carrier plus $n$ times the modulator frequency)
    \item Amplitude: $A \cdot J_n(\beta)$ (scaled by the $n$-th order Bessel function of $\beta$)
\end{itemize}
The Bessel functions thus directly determine how energy is distributed across the spectral components in an FM signal, explaining why they are fundamental to understanding and analyzing frequency modulation.
\begin{figure}[hbt!]
    \centering
    \includegraphics[width=0.95\linewidth]{images/desmosbessel_blank.png}
    \caption{Graph of Bessel functions of the first kind $J_n(x)$ for $n=\{0, 1, \dots, 5\}$, corresponding to red, orange, green, blue, purple, black respectively~\cite{desmos_bessel}}
    \label{fig:desmosbessel}
\end{figure}
\clearpage
\section{Amplitude Modulation (AM)}
Amplitude modulation (AM) represents another fundamental technique in signal processing that complements FM synthesis. While FM varies the frequency of a carrier signal, AM varies its amplitude according to a modulating signal.

\subsection{Basic Principles of AM}
In its simplest form, amplitude modulation multiplies a carrier signal by a modulator signal:
\begin{equation}
  AM(A, c, m) = A \cdot [1 + \mu \sin(mt)] \cdot \sin(ct)  
\end{equation}
Where:
\begin{itemize}
    \item $A$ is the carrier amplitude
    \item $c$ is the carrier frequency
    \item $m$ is the modulator frequency
    \item $\mu$ is the modulation depth (where $0 \leq \mu \leq 1$ for traditional AM)
\end{itemize}

The modulation depth $\mu$ controls how strongly the amplitude of the carrier is affected by the modulator. When $\mu = 0$, no modulation occurs; when $\mu = 1$, the carrier amplitude varies between $0$ and $2A$.
\begin{figure}[hbt!]
    \centering
    \includegraphics[width=0.85\linewidth]{images/am_waveform.png}
    \caption{AM waveform with increasing $\mu$}
    \label{fig:am_waveform}
\end{figure}
\subsection{Spectral Characteristics}
To understand the frequency domain representation of AM, we can expand the equation:

\begin{align}
AM(A, c, m) &= A \cdot [1 + \mu \sin(mt)] \cdot \sin(ct) \\
&= A \cdot \sin(ct) + A\mu \sin(mt)\sin(ct)
\end{align}

Using the trigonometric identity $\sin(\alpha)\sin(\beta) = \frac{1}{2}[\cos(\alpha-\beta) - \cos(\alpha+\beta)]$, we get:

\begin{align}
AM(A, c, m) &= A \cdot \sin(ct) + A\mu \cdot \frac{1}{2}[\cos(mt-ct) - \cos(mt+ct)] \\
&= A \cdot \sin(ct) + \frac{A\mu}{2}\cos((m-c)t) - \frac{A\mu}{2}\cos((m+c)t)
\end{align}

Converting to a uniform representation with sines:

\begin{align}
AM(A, c, m) &= A \cdot \sin(ct) + \frac{A\mu}{2}\sin((m-c)t + \frac{\pi}{2}) - \frac{A\mu}{2}\sin((m+c)t + \frac{\pi}{2}) \\
&= A \cdot \sin(ct) + \frac{A\mu}{2}\sin((c-m)t) + \frac{A\mu}{2}\sin((c+m)t)
\end{align}

This reveals the key characteristic of AM: the signal consists of the carrier frequency $c$ plus two symmetric sidebands at $c-m$ and $c+m$. Unlike FM which theoretically produces infinite sidebands, AM produces exactly three spectral components (assuming a single sine wave modulator):

\begin{itemize}
    \item The carrier at frequency $c$ with amplitude $A$
    \item The lower sideband at frequency $c-m$ with amplitude $\frac{A\mu}{2}$
    \item The upper sideband at frequency $c+m$ with amplitude $\frac{A\mu}{2}$
\end{itemize}

This spectral simplicity makes AM more predictable but less spectrally rich than FM for the same number of oscillators.

\begin{figure}[hbt!]
\centering
\includegraphics[width=0.9\linewidth]{images/am_spectrum_2.png}
\caption{Amplitude Modulation spectrum showing carrier and symmetric sidebands (Red: $\mu = 0$, Blue: $\mu = \frac{1}{2}$, Green: $\mu = 1$).}
\label{fig:am_spectrum}
\end{figure}


\chapter{The Frequency Modulation Transform}
\label{ch:6}
Although the technique of frequency modulation has been widely used across many signal processing disciplines for the last several decades, the exploration of more complex operator arrangements has been predominantly confined to musical instrument design. In this chapter, we propose an entirely new approach to signal analysis and decomposition leveraging the power of complex FM operator arrangements and their unique spectral characteristics. We hope to establish a new lens through which to view FM and create a set of tools to advance the domain of digital signal processing. Specifically, by the introduction of the Frequency Modulation Transform (FMT), which can decompose a signal into an FM arrangement with multiple modulators, each corresponding to that signal's most relevant harmonic intervals and their underlying associated modulation indices. As this type of analysis is algorithmic and purely deterministic, it allows for difficult signal processing tasks such as source identification and separation to be addressed without the use of deep learning. 
\section{Timbral Functions}
\subsection{Motivation}
Given most musical instruments (including human vocals) derive their timbre from the set of overtones (harmonic partials) produced when a sound is made, we can model the essence of a sound simply by looking at the arrangement of its partials and their amplitudes. Importantly, the partials produced by an instrument inherently depend on the fundamental frequency of the sound being produced. If we want to model this instrument mathematically for analysis or reconstruction, we must ensure that we capture the timbral characteristics as a function of the fundamental frequency. This connection lies at the heart of the technique of additive synthesis, in which a signal is broken down into its harmonics plus the fundamental, which are then re-synthesized to create a new instrument based on the spectral characteristics of the original sound~\cite{additive_2000}. However, we can take this further and analyze the spacings between partials to identify multiple distinct fundamentals for a multi-source audio and approximate the timbral characteristics for the underlying sources. 
\subsection{Timbral Functions}
Before we discuss how this decomposition works, we must first establish a framework for an instrument as a timbral function that produces a set of harmonics with amplitudes $A_k$ at multiples of the fundamental frequency $f$. For purposes of demonstration, we assume $k \in \mathbb{Z}$, but we can have non-integer (inharmonic) partials as well. 
\begin{definition}
    Let instrument $I:\mathbb{R} \to \mathbb{R}^h$ be a \textbf{timbral function} such that for some frequency $f \in \R$, $f\mapsto (A_1f,A_22f, \dots, A_hhf)$ where $h$ is a maximum harmonic that can be resolved for a given sampler. 
\end{definition}
In other words, for a signal $\bf{S}$ composed entirely of sounds from $I$, with fundamentals $f_1, f_2, \dots, f_m$ then the Discrete Fourier Transform $\mathcal{F}\{\bf{S}\}$ can be expressed as:
\begin{equation}
  \sum_{k=1}^h \sum_{j=1}^m A_ke^{-2\pi ikf_j} = \sum_{k=1}^hA_k\sum_{j=1}^m e^{-2\pi ikf_j}
\end{equation}
Notice that by construction, the amplitudes of the overtone series $\{A_k\}_{k=1}^\infty$ uniquely identify an instrument.\footnote{For real instruments this is infidelity a simplifying assumption as the spectrum can change depending on pitch/envelope, but for our purposes this assumption is sufficient} 
\newpage
\subsection{Frequency Modulation as a Timbral Function}
So far, we have only discussed simple FM operator arrays, which have one carrier and one modulator. Recall that the simple FM operator array with carrier $c$, modulator $m$, and modulation index $\beta$ has waveform $\sin(ct + \beta \sin(mt))$ and frequency spectrum:
\begin{equation}
  \sum_{n=-\infty}^\infty J_n(\beta)\sin((c+nm)t)  
\end{equation}
\begin{figure}[hbt!]
    \centering
    \includegraphics[width=0.6\linewidth]{images/simplefmfix.png}
    \caption{FM with a single carrier at 1200Hz and single modulator at 90Hz with modulation index $\beta= 0.5$ }
    \label{fig:simplefm}
\end{figure}

However, a common type of more complicated operator arrangement involves more than one modulator. For some carrier $c$ and modulators $m_1, m_2, \dots, m_k$ which have associated modulation indices $\beta_1, \beta_2, \dots, \beta_k$, the FM array has corresponding waveform:
\begin{equation}
  FM_{multi}(t) = \displaystyle\sin\left(ct + \sum_{i=1}^k \beta_i \sin(m_it)\right)  
\end{equation}
For which the resulting spectrum will take the form~\cite{Benson_2006}\footnote{\textit{Music: A Mathematical Offering} p 272}: 
\begin{equation}
  \sum_{n_1=-\infty}^\infty\sum_{n_2=-\infty}^\infty \dots \sum_{n_k= -\infty}^\infty \left[\displaystyle\Pi_{i=1}^k J_{n_i}(\beta_i)\right]\sin\left(c + \sum_{i=1}^kn_im_i\right)  
\end{equation}
Clearly, this is only slightly more complicated than the base case.\\
\begin{figure}[hbt!]
    \centering
    \includegraphics[width=0.6\linewidth]{images/multifm.png}
    \caption{FM with $c = 1200$Hz, three modulators $m_1 = 90$Hz, $m_2 = 72$Hz, $m_3 = 57$Hz with $\beta_1 = 0.5, \beta_2 = 0.5, \beta_3 = 0.1$ }
    \label{fig:miltifm}
\end{figure}
This new arrangement with multiple modulators is better for our purposes because we theoretically do not care about the carrier. If we are trying to find the set of shared harmonics of a set of frequencies, we can model each frequency as some modulator $m$ in an abstract FM arrangement. Moreover, the amplitudes on the overtone series can be in turn represented by the series of Bessel functions $\{|J_{k}(\beta)|\}_{k=0}^\infty$. This is useful because two frequencies with the same timbres will share a modulation index. More importantly, when we look at spacings between sidebands, we can model how harmonics from one instrument will interact with nearby harmonics from another. As the figure shows, despite originating from only three modulators, the interactions between FM operators can quickly become quite complicated.

Through this representation, we discovered a way to take a signal and obtain this set of modulators and modulation indices, which means that we can leverage properties and observations of this FM representation to analyze the original signals in ways that were not possible before.
\section{The Frequency Modulation Transform (FMT)}
\subsection{Motivation}
In order to find frequencies likely to be spaced at harmonic intervals, we must look at the distances between frequencies present in the spectrum of a signal. While this could be performed directly on the DFT through something like autocorrelation, these results can prove noisy, difficult to parse, and generalize poorly to more complex signals. To see why this is the case, refer back to figure \ref{fig:miltifm}. Looking at the DFT, we see a mixture of sidebands: some that correspond directly to integer multiples of the modulators, and others that are combinations of the modulators. For a real-world signal, this problem is amplified tremendously, as multiple sounds can occupy the same frequency ranges, which can obfuscate these harmonic intervals. Furthermore, as discussed in Section \ref{sec:leakage}, for signals with shorter sample lengths and thus wider DFT bins, determining the exact frequency spacings becomes even more difficult. Thus, we need an operation that takes a complicated signal and finds the underlying harmonics (modulators) in a way that is robust to the relationships between frequencies that are not harmonically related. In this section, we propose the Frequency Modulation Transform (FMT), which not only accomplishes the task of returning our desired intervals, but also yields the modulation indices that best encapsulate the relative amplitudes of each overtone and, therefore, encode its timbre.
\subsection{Mathematical Foundation}
Conceptually, the Fourier Transform can be understood as wrapping a signal around itself and measuring its periodicity at each frequency by the overlap between corresponding wrapped segments~\cite{3B1B}. While phase information is important for reconstruction, the magnitude of each component is precisely what tells us the strength of this periodicity. Hence, if we wanted to measure the periodicity of an instrument's underlying harmonics, we would expect large values at integer (and half integer) multiples of the fundamental, and theoretically zero response everywhere else. Thus, to find the patterns within a signal's frequency spectrum, or the ``frequencies of frequencies,'' we might think to simply just take the Fourier Transform again. 
\subsubsection{Breaking the Duality Property}
The issue with this approach is that taking the Fourier Transform more than once will just give us back a time-reversed version of the original signal. As for any periodic signal $x(t)$:
\begin{equation}
  \mathcal{F}\{\mathcal{F}\{\mathcal{F}\{\mathcal{F}\{x(t)\}\}\}\} = x(t)  
\end{equation}
The reason for this is that if the Fourier Transform is defined as a unitary operator:
\begin{align*}
    X(f)&=\mathcal{F}\{x(t)\}=\int_{-\infty}^\infty x(t)e^{-j2\pi f t}dt\\
    x(t)&=\mathcal{F}^{-1}\{X(f)\}=\int_{-\infty}^\infty X(f) e^{+j2\pi f t}df\\
\end{align*}
We can see that $+j$ and $-j$ are essentially interchangeable. From which we can observe the duality property~\cite{MIT_2011}:
\begin{equation}
    x(t) \xrightarrow[]{\mathcal{F}}X(f) \iff x(f) \xrightarrow[]{\mathcal{F}}X(t) 
\end{equation}
It follows that for even signals
\begin{align*}
    x(-f) &= \mathcal{F}\{X(t)\} = \int_{-\infty}^\infty X(t)e^{-j2\pi f t}dt\\
    X(-t) &= \mathcal{F}\{x(f)\} = \int_{-\infty}^\infty x(f)e^{+j2\pi f t}df \\
    \implies x(-t) &= \mathcal{F}\{\mathcal{F}\{x(t)\}\} \\
    \implies x(t) &= \mathcal{F}\{\mathcal{F}\{x(-t)\}\} \\
    \therefore x(t) &= \mathcal{F}\{\mathcal{F}\{\mathcal{F}\{\mathcal{F}\{x(t)\}\}\}\}
\end{align*}
Since in a DFT we assume a signal of length $N$ is periodic with period $N$, and for real-valued signals the DFT exhibits conjugate symmetry about the Nyquist frequency, the real parts of the components at each level will always be even. Therefore, successive applications of the Fourier Transform without modifications will not provide us with any new information beyond the duality property.

As a consequence, we will not be able to obtain the harmonic spacings through the Fourier transform alone. To accomplish this, we must design an operation that strategically breaks the duality property without destroying the harmonic information and allows us to go several levels deeper.

Notice that the assumptions of the duality property only hold if the signal is even and Hermitian, and, additionally, that the output of the first DFT is a sequence of complex (two-dimensional) vectors. Hence, if we want to take a second consecutive one-dimensional DFT, we have to reduce the dimensionality of our spectrum and break its symmetry. 

To accomplish both tasks, we define a new signal:
\begin{equation}
 \mathbf{S}' = \displaystyle\left\{||\mathcal{F}\{\mathbf{S}\}_i||_2\right\}_{i=1}^{N/2}   
\end{equation}
Where each element of $\mathbf{S}'$ is given by the magnitude of the corresponding complex component in $\mathcal{F}\{\mathbf{S}\}$. By only using the amplitudes, we know that all elements $\mathbf{s}_i' \in \mathbf{S}'$ will be non-negative, which prevents terms from destructively interfering when computing consecutive DFTs. Moreover, we ignore all components above the Nyquist frequency. Not only does this break symmetry, it prevents the second-order spectrum from containing frequencies that correspond to spacings across the Nyqusist frequency and across the zero frequency. To formalize the operation of symmetry-breaking, we can define it as follows:
\begin{definition}
    The \textbf{Breaking Operator}, $\mathcal{B}: \mathbb{C}^N \to \mathbb{R}_{\geq0}^{N/2}$, acts on a signal $\mathbf{S}$ of length $N$ such that $\mathbf{S} \mapsto \{||\mathbf{s}_n||_2\}_{n=1}^{N/2}$, where $\mathbf{s}_n\in \mathbb{C}$ is the $n$-th element of $\mathbf{S}$.
\end{definition}
Note that keeping up to $N/2$ frequencies is sufficient for the FMT, so successive applications of $\mathcal{B}$ will not reduce the lengths of the resulting spectra below $N/2$.

\subsubsection{Constructing the Transformation}
Recall that the purpose of the FMT is to find the ``frequencies of the frequencies.'' For clarity, we will disambiguate the two by explicitly defining these higher-order frequencies:
\begin{definition}
    A \textbf{k-th-order frequency} of a signal $\mathbf{S}$ is denoted by $f_i^{(k)} \in \mathbb{C}$ and is defined as the $i$-th element of the sequence given by the $k$-th consecutive application of a Discrete Fourier Transform to an asymmetric amplitude spectrum $\mathcal{F}\{\mathcal{B}\{\mathbf{S}^{(k - 1)}\}\}$.
\end{definition}
Note that a single sample in $\mathbf{S}$ is equivalent to a zero-th-order frequency, and components of the typical DFT can be thought of as first-order despite the lack of a breaking operation before the first DFT.

Hence, we can represent the harmonic relationships in a signal as higher-order frequencies, whose amplitudes correspond to the degree to which that relationship is expressed in the original signal. To go about extracting these higher-order frequencies, a single successive DFT will not be sufficient. To see why this is the case, recall that applying the DFT twice gives you a time-reflected version of the original signal. Even with the earlier construction of $\mathbf{S}'$, the resulting spectrum of $\mathcal{F}\{\mathbf{S}'\}$ still corresponds to a time-inverted version of the original source signal without all phase information and with strictly positive components. In an audio context, it is helpful to think of all even-ordered (zero is even) DFTs as producing a ``waveform,'' and odd-ordered DFTs as producing a ``frequency spectrum.'' 
\begin{figure}[hbt!]
    \centering
    \includegraphics[width=0.9\linewidth]{images/secondorderdft.png}
    \caption{Graph of second-order DFT amplitudes of FM array from Figure \ref{fig:simplefm} retains the wave-like characteristics of the original signal.}
    \label{fig:simplefmt}
\end{figure}
Hence, to obtain the desired spacings, we need frequencies of yet one order higher and must again break the symmetry of $\mathbf{S}'$ and take a third DFT. Now, we can formally define the FMT:
\begin{definition}
    The \textbf{Frequency Modulation Transform}, $\mathcal{M}:\mathbb{C}^N\to \C^{N/2}$ is defined as the mapping from a signal $\mathbf{S}$ of length $N$ to its corresponding sequence of third-order frequencies $\{f^{(3)}_i\}_{i=1}^{N/2}$, such that:
    \begin{equation}
    \mathcal{M}\{\mathbf{S}\} := \mathcal{F}\{\mathcal{B}\{  \mathcal{F}\{\mathcal{B}\{  \mathcal{F}\{\mathbf{S}\}  \} \}  \} \} 
\end{equation}
Where $\mathcal{F}$ denotes the Discrete Fourier Transform and $\mathcal{B}$ denotes the Breaking Operator.
\end{definition}

\subsection{Connections between FMT Amplitudes and FM Parameters}
Given the rather contrived definition of the FMT, the fact that it works at all is honestly quite surprising. Furthermore, the fact that it only works when breaking fundamental properties of the underlying signal is also rather strange and frankly somewhat unsettling. Nonetheless, the output of the FMT is extremely interesting, as the third-order frequencies correspond to the set of all possible harmonic intervals and their amplitudes correspond to the strength of that harmonic relationship (conceptually, this can be thought of as some measure of ``confidence'' in a third-order frequency corresponding to a modulator in some FM operator array, hence the name ``Frequency Modulation Transform''). 

All the more surprising is the relationship between the amplitudes on each third-order frequency and the underlying modulation indices with which they are associated. Initially, we expected these magnitudes to be somewhat abstract and altogether meaningless, but they are actually directly related to the modulation index and in such a way that the modulation index can be directly computed with a high degree of precision.

To illustrate the power of the FMT, let us return to the simple FM operator shown in Figure \ref{fig:simplefm} with carrier frequency $c = 1200$Hz, modulating frequency $m = 90$Hz and modulation index $\beta = 0.5$. This array has the following FMT:
\begin{figure}[hbt!]
\label{fig:simplefmt}
    \centering
    \includegraphics[width=0.675\linewidth]{images/fmtsimple.png}
    \caption{Normalized graph of the FMT of an FM operator with modulator $m = 90$Hz, produces large spike at $90$Hz, smaller at $2m = 180$Hz. }
\end{figure}
\clearpage
Here, we can see that the FMT is incredibly clean, with precisely three distinct peaks: the largest at the zero-frequency, the second-largest at the true modulating frequency, and the third at twice the modulating frequency. Furthermore, notice that the graph does not depend whatsoever on the carrier frequency and is identical if this value is changed to any frequency for which the highest sideband sits below the Nyquist frequency. 

More interesting is the FMT of the more complicated operator in Figure \ref{fig:miltifm} with three modulators: $m_1 = 90$Hz, $\beta_1 = 0.5$; $m_2 = 72$Hz, $\beta_2 = 0.5$; $m_2 = 57$Hz, $\beta_2 = 1.0$ 
\begin{figure}[hbt!]
    \centering
    \includegraphics[width=0.8\linewidth]{images/fmtmulti.png}
    \caption{Normalized graph of the FMT of an FM operator with three modulators has largest three spikes at the three modulating frequencies. }
    \label{fig:multifmt}
\end{figure}

Although this FMT is a bit messier than for the simple operator case, the largest non-zero peaks are still at the largest modulating frequencies. Moreover, the FMT amplitude is the same at $m_1$ and $m_2$, for which $\beta_1 = \beta_2$. The FMT amplitudes being identical for modulators with identical modulation indices shows that there is a fundamental connection between the FMT and FM operators with multiple modulators. 

If there is a way to go back and forth analytically between the FMT and its underlying signal, then we can create an FM operator that has theoretically identical FMT components. To do this, we must be able to decode FMT coefficients into their underlying modulation indices, which requires a useful observation about the structure of each of these graphs: FMT amplitudes correspond to underlying modulation index, but the zero frequency always has the highest magnitude. In the context of an FM array, a modulation index of zero simply produces the carrier; therefore, the amplitude of the third-order zero-frequency likely corresponds to the combined energy of the ``lonely'' lower-order frequencies that do not share harmonic relationships with any others. 

\subsection{Obtaining the Modulation Index}
Recall from Section \ref{sec:FM_synthesis} that increasing the modulation index ``steals'' energy from the carrier. Thus, the harmonic content in a signal should equivalently steal energy from the inharmonic content. Using this observation, we can begin to invert FMT amplitudes into their corresponding modulation indices.

If we plot the change in FMT amplitude for the largest non-zero third order frequency of a simple FM operator relative to its modulation index, we obtain the following set of points:
\begin{figure}[hbt!]
    \centering
    \includegraphics[width=0.65\linewidth]{images/points1alonefixed2.png}
    \caption{Measured FMT amplitude of largest-magnitude third-order frequency vs modulation index}
    \label{fig:pointsalone}
\end{figure}

Notice that for small values of $\beta$, this closely follows the shape of the first-order Bessel function $J_1(\beta)$, but then starts to diverge and remain above zero.

\begin{figure}[hbt!]
    \centering
    \includegraphics[width=0.7\linewidth]{images/points1greenline.png}
    \caption{Largest FMT amplitude vs modulation index with $J_1$}
    \label{fig:pointsline}
\end{figure}

Interestingly, the plotted points indicate that the underlying function appears to posses multiple sharp corners, which likely come from either squaring or taking the norm of DFT components in the preceding operations. Notable also is that there appears to be a sharp corner at the the peak of the first lobe ($\beta \approx 1.20241$) which is precisely the first zero of $J_0(2\beta)$. This is surprising, as we have not introduced any coefficients inside the Bessel calculations. Conceptually, however, an FM arrangement with modulator $m$ will have a sideband amplitude at frequency $c + nm$ of $J_n(\beta)$. Similarly, an arrangement with modulator $m_1 = nm \implies \sin(c+nm)$ will have a coefficient of $J_1(\beta)$, so there is some degree of equivalence between modulator coefficients when looking at different scales.\footnote{These nested coefficients on $\beta$ also appear when computing the frequency spectrum of FM arrangements with nested modulators (See \textit{Music: A Mathematical Offering} p. 273)} This hints that more is going on with these coefficients than it would seem simply from directly computing the DFT.  


% Interestingly, the FMT magnitudes for the second-highest amplitude frequency (''false'' modulator) similarly follows the second-order Bessel function $J_2(\beta)$:
% \begin{figure}[hbt!]
%     \centering
%     \includegraphics[width=0.6\linewidth]{images/points2blackline.png}
%     \caption{Second-largest FMT amplitude vs modulation index with $J_2$}
%     \label{fig:points2line}
% \end{figure}

While we were unfortunately unable to compute the explicit formula for the FMT amplitude of a true modulator given its modulation index index $\beta$ numerically, we believe with a high degree of certainty that it takes the form:
\begin{equation}
      \left|J_1(\beta)\right|\left[1 + J_0(\beta)^2\left(\sum_{n=1}^\infty J_n(n\beta)^2 \right) -  (1-J_0(\beta))^2\left(\prod_{n=1}^\infty(1-J_n(\beta)^2)\right) \right]
      \label{eq:bigdaddy}
\end{equation}
Although this equation may seem complicated, it comes as a more general summation taken from properties of the Bessel function combined with the formula for the frequency spectrum of an FM operator with multiple modulators. It is important to note that the largest normalized FMT coefficient in the single modulator case is equal to that same normalized coefficient in the multiple modulator case if all true modulators share the same modulation index. This is counterintuitive and it is unclear why this is the case. 

However, for all indices $\beta \in [0,1.20241]$, the function is one-to-one, so for the purpose of inverting the function to find the underlying modulation index, which will most likely lie in this region for practical FM arrangements, the approximation is sufficient. 
\clearpage
Figure \ref{fig:fmtmiapprox} displays the graph of the Equation \ref{eq:bigdaddy} and contains four curves:
\begin{align*}
    y_1 &=  \left|J_1(\beta)\right|\left[1 + |J_0(\beta)|^2\left(\sum_{n=1}^\infty |J_n(n\beta)|^2 \right) \right] \textbf{  (Blue)}\\
    y_2 &= \left|J_1(\beta)\right|(1-|J_0(\beta)|)^2\left(\prod_{n=1}^\infty(1-|J_n(\beta)|^2)\right) \textbf{  (Red)}\\
    y_3 &= y_1 - y_2 \textbf{  (Purple)} \\
    y_4 &= J_1(\beta)  \textbf{  (Green)}
\end{align*}
\begin{figure}[hbt!]
    \centering
    \includegraphics[width=1\linewidth]{images/fmtmiapprox_nogridlines.png}
    \caption{Function of largest FMT amplitude with respect to modulation index with Bessel functions evaluated up to order $n=6$.}
    \label{fig:fmtmiapprox}
\end{figure}
\newline
The error between the approximation and the peak points can be improved through attentional terms in $y_1$ of the form $(|J_n(n\beta)|^2)$, but because there is an $n$ inside the parenthesis, computing this accurately is difficult due to numerical stability issues with later terms in the power series expansion of $J_n(n\beta)$.
\clearpage
\subsection{Extending the FMT to Real-World Signals}
While FMT works well when dealing strictly with FM operators, it also generalizes to more complicated real-world signals. That said, harmonic relationships with real audio can be more complex than in the case of FM, and more often than not, have overtone coefficients that cannot be succinctly expressed with only one modulator. 

If we take the FMT of an audio segment with guitar, bass, banjo, and vocals, we can see that the peaks at which potential modulators exist are not as clearly separated from the ``chaff'' third-order frequencies as they were with FM.
\begin{figure}[hbt!]
    \centering
    \includegraphics[width=0.8\linewidth]{images/realworldfmt.png}
    \caption{FMT of real-world complex audio}
    \label{fig:realworldfmt}
\end{figure}
Not only is the zero frequency signal much more pronounced, there are a lot of peaks adjacent to others with similar magnitudes. However, the peaks at 49Hz, 147Hz, and those between 170Hz and 250Hz correspond to the fundamentals of the bass, guitar, and banjo respectively. Similarly, the FMT range for the vocals is a bit higher at 300-350Hz. The issue is distinguishing between third-order frequencies that actually correspond to harmonics.
\begin{figure}[hbt!]
    \centering
    \includegraphics[width=0.8\linewidth]{images/fmtpointspread.png}
    \caption{Normalized amplitude on 100 largest FMT peaks}
    \label{fig:fmtpointspread}
\end{figure}
Also, as lower frequencies tend to have higher magnitudes at equal loudness to higher frequencies, the lower FMT components can appear over expressed compared to their perceptual importance.

Another interesting observation is the presence of low-frequencies present in the FMT. These reflect the character of peaks that are not isolated to one bin (or between two, which can also explain why some many FMT peaks are close together). Given this is one second audio, the bins will be relatively wide, which increases the chance that a frequency will fall closer to the edge of a bin than the middle. 

The largest nonzero FMT peaks corresponds 49Hz,  which agrees with the fundamental of the known bass note being played during the audio. Interestingly, if we look at the initial FFT of this audio in Figure \ref{fig:fftof1s}, we can see that 49Hz (G1) has a much lower amplitude than its first harmonic at 98Hz (G2), which is the actual pitch of the note played on the bass. The second highest FMT amplitude sits between to 24 and 25Hz, which corresponds to (G0), which is not directly present in the FFT but is the ``true'' fundamental frequency for the note G.
\clearpage
\begin{figure}[hbt!]
    \centering
    \includegraphics[width=0.8\linewidth]{images/fftof1s.png}
    \caption{FFT of real world audio}
    \label{fig:fftof1s}
\end{figure}
We can also see three spikes at 98Hz (G2), 123Hz (B2), and 147Hz (D3), which correspond to the fundamental frequencies of the notes comprising the G major chord, which the guitar is playing during the recording. The peak at G2 is larger due to the additional frequencies of the bass. Additionally, peaks at 196Hz, 250Hz, and 294Hz correspond to the subsequent harmonics of each of the notes in chord, and we see these same intervals expressed in the FMT.

\clearpage
\subsubsection{FMT Reconstruction with FM}
If we take the largest 10 peaks and solve for the associated modulation index, we can construct an FM arrangement that should theoretically have a similar FMT.

What we see in Figure \ref{fig:combinedfmt} is that even with 10 modulators and an approximation of their modulation indices, the two have strikingly similar FMT graphs, with nearly identical amplitudes at the chosen peaks.
\begin{figure}[hbt!]
    \centering
    \includegraphics[width=0.8\linewidth]{images/combinedfmt.png}
    \caption{Normalized FMT of real-world audio (blue), Normalized FMT of FM approximation with 10 modulators(red)}
    \label{fig:combinedfmt}
\end{figure}

Interestingly, the red (FM FMT) graph looks identical for every reasonable\footnote{Here, ``reasonable'' implies that the carrier is nonzero and far enough away from the Nyquist frequency that sidebands do not extend past this limit.}  choice of carrier.

The connection between FM synthesis and real-world acoustic instruments is not coincidental. In the 1980s, Yamaha engineers discovered that carefully configured FM operator arrangements could convincingly emulate pianos, brass, strings, and other acoustic instruments~\cite{DX7}. This breakthrough with the DX7 synthesizer revealed that the Bessel function relationships inherent in FM synthesis naturally mirror the harmonic structures found in many acoustic timbres, particularly those with dynamic spectral content. The success of FM in reproducing diverse acoustic phenomena stems from its ability to model time-varying spectral characteristics—just as the overtone amplitudes of real instruments change during attack, sustain, and decay phases, FM's modulation index can dynamically vary the distribution of energy across its sidebands. This mathematical alignment explains why the FMT can extract meaningful harmonic relationships from real-world audio, even though natural instruments produce more complex spectra than pure FM operators. The fact that FM synthesis can approximate such a wide range of acoustic timbres suggests that Bessel function-based spectral decomposition represents a fundamental way of understanding and modeling natural sound.

Ultimately, we found these initial results to be extremely exciting, which leads us to believe there is much more to be done in this area. The FMT illuminates a previously-unexplored area of signal processing, which opens the door for more nuanced timbral analysis and a new tool to use in audio processing.
\newpage

\chapter{FMT: Applications}
\label{ch:7}
\section{Signal Decomposition and Analysis}
Given we can now represent the harmonics and timbral information of sound using a (comparatively) simpler FM representation that mimics its structure, we can leverage properties of this FM arrangement to give insight into which base frequencies belong to a given timbral collection. Specifically, we can try to ``invert'' parts of the FMT using several key observations about the DFT: 
\begin{enumerate}
    \item By keeping $N/2$ samples, $\mathcal{B}$ is not destructive to the magnitudes of the first DFT. In other words, we can always infer the ``missing'' half of the DFT through Hermitian symmetry.
    \item If we retain a copy of each successive DFT before applying  $\mathcal{B}$, we can simply replace the old magnitudes with new ones obtained from inverse transforms and use the copied phases.
    \item Although the carrier in the corresponding FM representation does not matter when approximating the FMT of the signal, choosing carriers corresponding to high magnitude frequencies in the original spectrum and comparing their DFTs can yield insight into which frequencies might correspond to which harmonics and vice versa.
\end{enumerate}
To better illustrate the third point, consider the DFT of the FM representation with carrier $c = 98$Hz, which corresponds to G2 and has the highest magnitude out of any frequency in the original DFT. 
\begin{figure}[hbt!]
    \centering
    \includegraphics[width=0.85\linewidth]{images/dftcomparisonfmt.png}
    \caption{Normalized FFT amplitudes of real-world audio (blue), Normalized FFT amplitudes of FM approximation with $c = 98$Hz (red)}
    \label{fig:dftcomparisonfmt}
\end{figure}
\newline
The important takeaway this graph is that informed choices of carrier frequencies will lead to an FM representation that captures both the structural characteristics of the FMT and the FFT. 

% In fact, if we subtract off part of the carrier, which which will always have a relatively high amplitude, with a sine 98Hz wave that is in phase we get a new graph that retains even more of these characteristics:
% \begin{figure}[hbt!]
%     \centering
%     \includegraphics[width=0.65\linewidth]{images/betterdftcomparison.png}
%     \caption{Normalized FFT amplitudes of real-world audio (blue), Normalized FFT amplitudes of FM approximation with $c = 98$Hz minus a sine wave at 98.0 (red)}
%     \label{fig:betterdftcomparison}
% \end{figure}


\section{Source Separation Without Machine Learning}
\subsection{Motivation}
As we have shown thus far, there is an intrinsic connection between the FMT representation of a signal and its original DFT. Moreover, this connection can be exploited to glean previously-unavailable contextual information about the signal's important frequencies and their relationships to one another.

The audio processing task of source separation, which breaks apart an audio composed of sounds from multiple sources (e.g. instruments, speakers, etc.) has been predominantly constrained to the domain of deep learning, as this task often involves nonlinear, contextual relationships hidden within frequency/time-frequency representations of the audio~\cite{ANSARI2023126895}. Even with sophisticated deep learning architectures, source separation is by no means trivial. At a fundamental level, these networks must learn to distinguish between instruments based on their timbres, which can overlap at important frequencies and ultimately confuse the model. However, with adequate domain knowledge coupled with the FMT, we can perform an approximation of source extraction with virtually zero aliasing and a drastically smaller compute requirement. 
\subsection{Algorithm - FMTSS}
The FMT Source Separation (FMTSS) algorithm is far from perfect and is intended to serve as an initial proof of concept and to reinforce the validity of using an FMT-based approach. That said, the algorithm is fairly straightforward:
\begin{enumerate}
    \item Take the FFT of an audio signal and store it. Then, take the FMT of that audio and construct an approximation of the underlying FM arrangement by inverting the FMT amplitudes to obtain the corresponding modulation indices (a dictionary with linear interpolation is used to circumvent expensive Bessel function calculations). 
    \item Either:
    \begin{itemize}
        \item Start with a prominent frequency belonging to the target source \\
        OR
        \item Start with any prominent frequency that is an integer multiple of a third-order frequency with a large coefficient. Importantly, this third-order frequency should be audible ($>$20-30Hz). 
    \end{itemize}
    \item (Optional) If the audio task involves instruments, analyze the DFT for information at octave intervals above the selected frequency (e.g. $2f$). 
    \item Identify the highest magnitude third-order frequency $f_k^{(3)} = \frac{1}{k}f : k \in \N$, where $f_k^{(3)} > 20$Hz, construct the simple FM operator: 
    \begin{equation}
      FM(f,f_k^{(3)}, \beta_k )(t) = \sin(2\pi f t+\beta_k\sin(2\pi f_k^{(3)}t))  
    \end{equation}
    \item Let $A_f \in \mathbb{R}_{\geq 0}$ be the amplitude of the DFT of the mixture signal at frequency $f$, let $\gamma \in [0.5, 1]$\footnote{$\gamma \approx 0.8$ or something in that ballpark works well.} be a damping term determined experimentally (or computed through something like LLS). Now, take $\mathcal{F}\{FM(f,f_k^{(3)}, \beta_k )(t)\}$. Let $B_f = |J_0(\beta_k)|$  be the magnitude of the FM frequency spectrum at carrier\footnote{If using multiple modulators just use the magnitude at $f$} $c =f$. Next, normalize and rescale its amplitude spectrum such that its carrier has amplitude $\gamma A_f$. If an additional octave frequency is being used, do this for $\gamma A_{2f}FM(2f,f_k^{(3)}, \beta_k )(t)$ as well. Mathematically, we can express this as:
    \begin{equation}
      \frac{\gamma A_f}{B_f }\mathcal{F}\{FM(f,f_k^{(3)}, \beta_k )(t)\}  
    \end{equation}
    \item (Optional) Add in sub-band (0-20Hz) third-order frequencies with high amplitudes to the FM operator(s). If it is prominent in the FMT of the signal, the 0.5-1Hz modulator can help to preserve details about the volume envelope of the original source in the reconstructed audio. 
    \item Swap the signal's original magnitudes with the rescaled FM magnitudes and take the IFFT using the original phases. 
\end{enumerate}


\subsection{Results}
Fine tuning the choice of low frequency modulators can substantially improve the qualitative output, but further steps could also be taken to automate this given the usage specification. That said, for what amounts to essentially five FFTs, one IFFT, and scalar multiplication, the results are extremely promising, especially considering the potentially exorbitant computational cost of using deep learning-based approaches which were previously the only option.
\clearpage
For the one second audio illustrated in the previous section, we set to isolate the electric bass with optimal FM parameters:
$$f = 98.0\text{Hz}, \;\{(m,\beta)\}  = \{(49.0\text{Hz},0.2725972), \,(1.0\text{Hz},686119), \,(2.0\text{Hz},0.396291)\}$$
And an additional octave operator with parameters:
$$2f = 196.0\text{Hz}, \;\{(m,\beta)\}  = \{(49.0\text{Hz},0.2725972), \,(1.0\text{Hz},686119)\}$$
\begin{figure}[hbt!]
    \centering
    \includegraphics[width=0.8\linewidth]{images/bassreconfft.png}
    \caption{Ground-truth electric bass stem (purple), reconstructed (blue), mixture audio (green)}
    \label{fig:bassreconfft}
\end{figure}
\newline
Not only does the reconstruction show frequency peaks at every single harmonic of the ground truth bass, our reconstruction is able to accomplish this with two relatively simple FM arrangements that only use modulators within the top ten third-order frequencies by magnitude. Importantly, Figure \ref{fig:bassreconfft} illustrates how we only need to solve for the peak amplitude for a given FM array, as the relative amplitudes on the harmonics (and fundamental) are already given by the modulation indices we obtained earlier. This saves a substantial amount of computation and yields fundamentally novel results, as the underlying timbral structure of the bass is baked-in to the FM representation. 

Worth noting is that there is a bit of error in approximating the exact frequency and amplitude of these harmonics, which stems mainly from the limitations of using audio that is only one second long (i.e. we have 1Hz bin resolution at 48kHz). Nonetheless, the fact that the approximation is this good even with relatively coarse resolution is still extremely promising. Moreover, this is without using DKD for more precise frequency estimation. 

Copying over the phases while only manipulating DFT magnitudes is a technique commonly used in DSP, as phase information is notoriously fragile. Even small phasing issues can lead to audible distortion. With our approach, we can take advantage of the initial phases, which are already known to be coherent without running the risk of bleed. This is because the FM signals only have nonzero amplitudes at frequencies within the span of the modulators. 
\begin{figure}[hbt!]
    \centering
    \includegraphics[width=0.8\linewidth]{images/bass_stft.png}
    \caption{STFT of reconstructed stem (top); ground truth stem (bottom). \\
    $\texttt{n\_fft} = 2048, \;\texttt{hop\_len}$ = 128. 1 time unit = 250ms}
    \label{fig:bassreconfft}
\end{figure}
As we can see from the above STFT, we can construct a convincing reconstruction by retaining the original phase information before the IFFT. Two important observations about the graph: firstly, the breakup after t=3.5 (875ms) in the ground truth is from artifacts caused by the author taking their finger off the fret prematurely and has nothing to do with the underlying harmonic structure or (desired) volume envelope; also the FMTSS model fails to capture the high frequencies present during the decay stage of the volume envelope (t=0.1-2.0). While additional modulators could fix this, dealing with more complex time-frequency changes in harmonics goes beyond the scope of this bare-bones initial implementation. 

Qualitatively, the two audio segments sound incredibly similar, and the reconstruction successfully captures the timbral character of the original despite lacking some nuances present in the envelope. Regardless, this proof of concept illustrates the promising nature of the FMT and should serve as a starting point for how future applications might leverage the FMT.  
\section{Future Work and Other Applications of the FMT}
Admittedly, the FMT and its applications were not as thoroughly explored in this paper as we would have liked.\footnote{The reason for this is that the connection between third-order frequencies and FM operators was discovered less than three weeks before the initial thesis deadline, and its connection to modulation indices several days after that.} While we plan to continue developing the FMT and making it more robust to time-frequency changes, we still felt it was important to include our initial findings in this paper. In this section we will discuss several other areas for improvement and potential applications for the FMT, a potential extension of the FMT to 2D, and finally, we end with a conjecture on the universality of FMT representations of acoustic signals. 
\subsection{Connection to DKD}
\label{subsec:dkdfmt}
An immediate connection can be drawn to Dirichlet Kernel Deconvolution as a way to improve the resolution at which we can determine possible modulators and carriers without increasing the sample rate or signal length. Ideally, we perform DKD both within the FMT and on the initial DFT itself, but given the breaking operators and compounding effects of adding successive kernels, DKD might require modifications to work with third-order frequencies directly. 
\subsection{Applications to Machine Learning}
Given our initial success with using the FMT to isolate a single source audio from a mixture, we feel that FMT might be well-suited for other applications previously restricted to machine learning. Moreover, FMT is not intended to be a replacement for more sophisticated techniques; rather, it is intended to work in tandem with existing approaches and can offer a new method through which to derive and inform frequency-domain features. 

\subsubsection{Classification Tasks}
Another audio processing task that primarily relies on deep learning techniques is source classification~\cite{5678725}. For example, distinguishing between different musical instruments or identifying speakers in an audio recording. The FMT provides a compact representation of harmonic relationships that could serve as distinctive features for classification models.
Traditional machine learning approaches for audio classification often rely on handcrafted features like MFCCs (Mel-frequency cepstral coefficients), spectral centroid, and spectral rolloff~\cite{labelstudio_2022}. These features, while useful, do not directly capture the harmonic structure of sound sources. The FMT, by encoding information about the fundamental frequencies and their relationships through modulation indices, offers a more semantically meaningful representation.

Specifically, a classification pipeline incorporating FMT might:
\begin{enumerate}
    \item Extract FMT features from audio samples
    \item Identify the most prominent third-order frequencies
    \item Calculate corresponding modulation indices
    \item Use these as feature vectors for traditional machine learning algorithms
\end{enumerate}
This approach could be particularly effective for classifying sounds with distinctive harmonic structures, such as musical instruments or speech from different speakers. Moreover, the dimensionality reduction inherent in the FMT representation could potentially improve training efficiency compared to approaches that use raw spectrograms or waveforms.
\subsubsection{Feature Engineering for Deep Learning}
Beyond traditional machine learning, the FMT could serve as a complementary input to deep neural networks. Recent advances in audio machine learning have emphasized end-to-end approaches where networks learn directly from minimally processed inputs. However, providing structured harmonic information through FMT features could help networks converge more quickly with improved generalization.

For example, a hybrid approach might:
\begin{enumerate}
    \item Process audio through both spectral and FMT pathways
    \item Feed the resultant features into a neural network
    \item Allow the network to learn which representation provides more discriminative information for different tasks
\end{enumerate}
This multi-pathway architecture could be particularly valuable for complex tasks like music transcription, where both time-frequency content and harmonic structure are important.
\subsubsection{Unsupervised Learning and Clustering}
The FMT's ability to extract harmonic relationships suggests applications in unsupervised learning contexts. Audio segments could be clustered based on their third-order frequency distributions, potentially revealing natural groupings of sound sources without labeled training data.

This could be valuable for tasks like audio database organization, where manual labeling is impractical. Moreover, the interpretability of FMT features—directly tied to harmonic relationships—makes the resulting clusters more meaningful and easier to analyze than those derived from black-box representations.
\clearpage
\subsubsection{Conclusion and Future Direction}
While our exploration of the FMT's applications to machine learning is preliminary, the results suggest several promising directions. The mathematical connection between third-order frequencies and FM parameters provides a bridge between signal processing theory and practical audio applications that has been largely untouched.
Future work might explore integrating FMT with attention mechanisms in neural networks, allowing models to focus specifically on the most relevant harmonic relationships for a given task. Additionally, extending the time-frequency resolution of FMT through techniques like wavelet transforms could help capture more dynamic aspects of audio signals.
The FMT represents not just a new analysis technique, but potentially a new paradigm for understanding audio signals through their harmonic structures—a paradigm that could complement and enhance existing machine learning approaches while providing greater interpretability and efficiency.
\subsection{Hypothetical Connections to the Polar Fourier Transform}
So far in this paper we have only discussed applications for one-dimensional signals; however, two-dimensional Fourier transforms are widely used in other DSP domains such as image processing. Despite images lacking the same notion of timbre that we might find in an audio signal, there is an intriguing connection between the structure of an FM operator's frequency representation and the basis functions for the polar Fourier transform in two dimensions. 

To explore this further, let us examine the structure of the radial basis function for a Polar Fourier transform defined on the whole space derived from radial solutions to the Helmholtz Equation~\cite{WRB08}: 
\begin{equation}
    \Psi_{k,m}(r,\varphi) = \sqrt{k}J_m(kr)\Phi_m(\varphi)
\end{equation}
Where $k\in \R_{\geq0}$ and $\Phi_m(\varphi)$ is defined by
\begin{equation}
    \Phi_m(\varphi) = \frac{1}{\sqrt{2\pi}}e^{im\varphi}
\end{equation}
However, this function is typically only defined on a finite region $r \leq a$ rather than the whole space. In this case a function can be expanded as:
\begin{equation*}
    f(r, \varphi) = \sum_{n=1}^\infty \sum_{m=-\infty}^\infty P_{nm}\Psi_{nm}(r,\varphi)
\end{equation*}
Where polar Fourier coefficient $P_{nm}$ is given by
\begin{equation}
    \int_0^a\int_{0}^{2\pi}f(r,\varphi)\Psi^*_{nm}(r, \varphi)rdrd\varphi
\end{equation}
Despite the author's lack of expertise in this specific area, the structure of the polar Fourier transform is interesting for two reasons: Firstly, the 
$$\sum_{m=-\infty}^\infty P_{nm}\Psi_{nm}(r,\varphi)$$
segment of the transform can be interpreted as a sum of frequencies with coefficients of Bessel functions of the first kind with the same bounds of summation as with the one dimensional Cartesian Fourier transform for FM operators discussed in Section \ref{sec:2_2}. Second, the values of $k$ relate to the energy levels of the system and their wave functions. So in the context of images, $k$ reflects the scale of basic patterns. This could indicate that there is a direct relationship between classes of ``polar modulation indices'' in a hypothetical polar FM operator and radial basis functions. If this connection can be expressed in terms of some higher-dimensional, polar analogue to the FMT, then the operation might be even better suited to 2D signals than 1D.

Admittedly this all goes well beyond both the scope of the author's own understanding as well as the scope of this paper; however, we felt that the similarities in the structure of FM sidebands and the Polar Fourier transform were worth mentioning--especially given our initial success with the FMT in one dimension. 
\clearpage
\section{Conjecture: Universality of FMT Decomposition}
\begin{conj}[Universality of FMT Decomposition]\label{conj:universality}
For any periodic acoustic signal $s(t)$ with finite harmonic content, there exists a finite-dimensional frequency modulation operator arrangement $\mathbf{M}^* = \{(c_i, m_i, \beta_i)\}_{i=1}^N$ such that:
\begin{enumerate}
    \item The Frequency Modulation Transform coefficients $\mathcal{M}\{s(t)\}$ uniquely determine the signal's timbral characteristics up to perceptual distinguishability.
    
    \item The reconstruction error $\epsilon = \|s(t) - \tilde{s}(t)\|_2$ falls below the threshold of human auditory perception, where $\tilde{s}(t)$ is the signal reconstructed from $\mathbf{M}^*$.
    
    \item The number of operators $N$ is bounded by a function of the signal's spectral complexity, specifically $N \leq \mathcal{O}(\log H)$ where $H$ is the number of significant harmonics in $s(t)$.
\end{enumerate}
\end{conj}

This conjecture suggests that frequency modulation synthesis, through its intrinsic connection to Bessel functions, provides a universal framework for representing acoustic timbre. If true, this would establish the FMT as a fundamental tool for audio analysis and synthesis, analogous to how the Fourier transform serves as a universal tool for frequency analysis.

The implications of this conjecture extend beyond signal processing to our understanding of how natural sounds are structured and perceived. The success of FM synthesis in emulating diverse acoustic instruments suggests that this mathematical framework may reflect fundamental principles of sound production in physical systems.
\addcontentsline{toc}{chapter}{\protect\numberline{}{Appendix}}
\chapter*{Appendix}
\section*{Nyquist-Shannon Sampling Theorem}
\label{sec:NQST}

The Nyquist-Shannon sampling theorem~\cite{Shannon1949} is a fundamental principle in digital signal processing that establishes the necessary conditions for accurately sampling and reconstructing continuous signals. Formally stated, the theorem asserts:

\textbf{Theorem} (Nyquist-Shannon Sampling Theorem). \textit{If a function $x(t)$ contains no frequencies higher than $B$ hertz, it is completely determined by giving its ordinates at a series of points spaced $\frac{1}{2B}$ seconds apart.} \\


More practically, this means that to perfectly reconstruct a bandlimited continuous-time signal from its samples, the sampling rate $f_s$ must be greater than twice the maximum frequency $B$ present in the signal:
$$f_s > 2B$$

The quantity $2B$ is commonly referred to as the Nyquist rate, and half the sampling rate ($\frac{f_s}{2}$) is known as the Nyquist frequency.

\subsection*{Implications for Digital Audio}

For digital audio, this theorem has profound implications:

\begin{itemize}
    \item \textbf{Sampling Rate Selection}: Since human hearing typically extends to about 20 kHz, digital audio is commonly sampled at 44.1 kHz or 48 kHz to satisfy the Nyquist criterion.
    
    \item \textbf{Anti-aliasing Filters}: Before sampling, analog signals must be filtered to remove frequencies above the Nyquist frequency to prevent aliasing.
    
    \item \textbf{Aliasing}: If a signal contains frequencies above the Nyquist frequency, these will ``fold back'' and appear as lower frequencies in the sampled signal, creating distortion.
\end{itemize}

\subsection*{Mathematical Formulation}

The reconstruction of a continuous-time signal $x(t)$ from its samples $x[n] = x(nT)$ (where $T = \frac{1}{f_s}$ is the sampling period) can be expressed as:

$$x(t) = \sum_{n=-\infty}^{\infty} x[n] \cdot \text{sinc}\left(\frac{t-nT}{T}\right)$$

where $\text{sinc}(x) = \frac{\sin(\pi x)}{\pi x}$ is the normalized sinc function.

This reconstruction formula, known as the Whittaker–Shannon interpolation formula, provides the theoretical basis for perfect reconstruction of a bandlimited signal from its samples, provided the Nyquist criterion is met~\cite{pollock_whittaker}.

\subsection*{Practical Considerations}

In practice, perfect reconstruction is rarely achieved due to:

\begin{itemize}
    \item \textbf{Finite-length signals}: Real-world signals have finite duration rather than extending infinitely.
    
    \item \textbf{Truncated sinc functions}: Practical reconstruction uses windowed or approximated sinc functions.
    
    \item \textbf{Quantization error}: Digital samples have finite precision, introducing quantization noise.
    
    \item \textbf{Non-ideal filters}: Practical anti-aliasing filters cannot perfectly eliminate all frequencies above the Nyquist frequency.
\end{itemize}

Despite these practical limitations, the Nyquist-Shannon sampling theorem provides the theoretical foundation for all digital signal processing and forms the basis for understanding the relationship between continuous analog signals and their discrete digital representations.

\section*{FM Operator Composition and Composite Operators}
\label{sec:2_2}
\begin{lem} Adding two operator arrays with opposite-signed modulating frequencies removes odd harmonics and doubles the amplitudes of even harmonics.
$$\displaystyle FM(A, c, m, \beta) + FM(A, c, -m, \beta) = 2A\sum_{n=-\infty}^\infty J_{2n}(\beta)\sin((c + 2nm)t)$$
\end{lem}

\begin{proof} First, we can write that by linearity of the Fourier transform, we have:
\begin{align*}
    FM(A, c, m, \beta) + FM(A, c, -m, \beta) &= \mathcal{F}(FM(A, c, m, \beta) + FM(A, c, -m, \beta)) \\
    &= \mathcal{F}(FM(A, c, m, \beta)) + \mathcal{F}(FM(A, c, -m, \beta))\\
    &=\Tilde{FM}(A, c, m, \beta) + \Tilde{FM}(A, c, -m, \beta)
\end{align*}
By definition of the Fourier transform of $FM$, we can expand this as:
\begin{align*}
    \Tilde{FM}(A, c, m, \beta) + \Tilde{FM}(A, c, -m, \beta) &= A \sum_{n=-\infty}^\infty J_n(\beta) \sin((c + nm)t) + A \sum_{n=-\infty}^\infty J_n(\beta) \sin((c + n(-m))t)\\
    &= A \sum_{n=-\infty}^\infty J_n(\beta) \sin((c + nm)t) + A \sum_{n=-\infty}^\infty J_n(\beta) \sin((c - nm)t)
\end{align*}
Now, note that for even $n$, $J_n$ is even $\implies J_{2n}(-\beta) = J_{2n}(\beta)$ and for odd $n$ we have that $J_{n}$ is also odd $\implies J_{n}(-\beta) = -J_{n}(\beta)$.\\

Now, we will group the two sums into infinitely many sums of the $n$-th and $-n$-th indices from both $\tilde{FM}(A, c, m, \beta)$ and $\tilde{FM}(A, c, -m, \beta)$ as follows:
\begin{align*}
    A \sum_{n=-\infty}^\infty J_n(\beta) \sin((c + nm)t) + A \sum_{n=-\infty}^\infty J_n(\beta) \sin((c - nm)t) \\
    = 2A J_0(\beta)\sin(ct) + A\sum_{n=1}^\infty [ J_{-n}(\beta)\sin((c-nm)t) + J_{n}(\beta)\sin((c+nm)t)  \\
    + J_{-n}(\beta)\sin((c-n(-m))t) + J_{n}(\beta)\sin((c+n(-m))t) ]
\end{align*}
We now have two cases:\\
For even $n = 2n$, we have:
\begin{align*}
    2A J_0(\beta)\sin(ct) + A\sum_{n=1}^\infty [ J_{-2n}(\beta)\sin((c-2nm)t) + J_{2n}(\beta)\sin((c+2nm)t)  \\
    + J_{-2n}(\beta)\sin((c-2n(-m))t) + J_{2n}(\beta)\sin((c+2n(-m))t) ]\\
    = 2A J_0(\beta)\sin(ct) + A\sum_{n=1}^\infty [ J_{2n}(\beta)\sin((c-2nm)t) + J_{2n}(\beta)\sin((c+2nm)t)  \\
    + J_{2n}(\beta)\sin((c+2nm)t) + J_{2n}(\beta)\sin((c-2nm)t) ] \\
    = 2A J_0(\beta)\sin(ct) + A\sum_{n=1}^\infty J_{2n}(\beta)[ \sin((c-2nm)t) \\
    + \sin((c+2nm)t) + \sin((c+2nm)t) + \sin((c-2nm)t) ]\\
    = 2A J_0(\beta)\sin(ct) + A\sum_{n=1}^\infty J_{2n}(\beta)[ 2\sin((c-2nm)t) + 2\sin((c+2nm)t) ]\\
    = 2A J_0(\beta)\sin(ct) + 2A\sum_{n=1}^\infty J_{2n}(\beta)\sin((c-2nm)t) + 2A\sum_{n=1}^\infty J_{2n}(\beta)\sin((c+2nm)t)\\
    = 2A \sum_{n=-\infty}^\infty J_{2n}(\beta)\sin((c + 2nm)t)
\end{align*}
However, for odd $n$, we have:
\begin{align*}
    A\sum_{n=1, 2\nmid n}^\infty [ J_{-n}(\beta)\sin((c-nm)t) + J_{n}(\beta)\sin((c+nm)t)  \\
    + J_{-n}(\beta)\sin((c-n(-m))t) + J_{n}(\beta)\sin((c+n(-m))t) ]\\
    = A\sum_{n=1, 2\nmid n}^\infty [ -J_{n}(\beta)\sin((c-nm)t) + J_{n}(\beta)\sin((c+nm)t)  \\
    + -J_{n}(\beta)\sin((c+nm)t) + J_{n}(\beta)\sin((c-nm)t) ] \\
    = A\sum_{n=1, 2\nmid n}^\infty [ \sin((c-nm)t)(J_{n}(\beta)-J_{n}(\beta)) + \sin((c+nm)t)(J_{n}(\beta) - J_{n}(\beta) ) \\
    = A\sum_{n=1, 2\nmid n}^\infty [ \sin((c-nm)t)(0) + \sin((c+nm)t)0 ) \\
    = A\sum_{n=1, 2\nmid n}^\infty 0 + 0 \\
    = 0
\end{align*}
Combining both odd and even $n$, we get:
\begin{align*}
\Tilde{FM}(A, c, m, \beta) + \Tilde{FM}(A, c, -m, \beta) &= 2A J_0(\beta)\sin(ct) \\
&+ A\sum_{n=1}^\infty [ J_{-n}(\beta)\sin((c-nm)t) + J_{n}(\beta)\sin((c+nm)t)  \\
    &+ J_{-n}(\beta)\sin((c-n(-m))t) + J_{n}(\beta)\sin((c+n(-m))t) ]\\
    &= 2A \sum_{n=-\infty}^\infty J_{2n}(\beta)\sin((c + 2nm)t) + 0 \\
    &= 2A \sum_{n=-\infty}^\infty J_{2n}(\beta)\sin((c + 2nm)t) 
\end{align*}
\end{proof}

\begin{lem} Adding two operator arrays with opposite-signed carrier frequencies removes even harmonics and doubles the amplitudes of odd harmonics.
$$\displaystyle FM(A, c, m, \beta) + FM(A, -c, m, \beta) = 2A\sum_{n=-\infty, 2\nmid n}^\infty J_{n}(\beta)\sin((c + nm)t)$$
\end{lem}


\begin{proof} This lemma follows from the first if we consider that because sine is odd:
$$\sin(((-c) + nm)t) = \sin((nm -c)t) = \sin(-1(c - nm)) = -\sin(c - nm)$$
$$= -\sin(c + n(-m))$$
Hence, for even $n$, we have:
\begin{align*}
    J_{-2n}(\beta)\sin((c-2nm)t) + J_{2n}(\beta)\sin((c+2nm)t) \\
    + J_{-2n}(\beta)\sin((-c-2nm)t) + J_{2n}(\beta)\sin((-c+2nm)t) \\
     = J_{2n}(\beta)\sin((c-2nm)t) + J_{2n}(\beta)\sin((c+2nm)t) \\
    + J_{2n}(\beta)\sin(-1(c+2nm)t) + J_{2n}(\beta)\sin(-1(c-2nm)t) \\
    = J_{2n}(\beta)\sin((c-2nm)t) + J_{2n}(\beta)\sin((c+2nm)t) \\
    - J_{2n}(\beta)\sin((c+2nm)t) - J_{2n}(\beta)\sin((c-2nm)t) \\
    =\sin((c+2nm)t)(J_{2n}(\beta) - J_{2n}(\beta)) + \sin((c-2nm)t)(J_{2n}(\beta)- J_{2n}(\beta)) \\
    =\sin((c+2nm)t)(0) + \sin((c-2nm)t)(0) \\
    = 0
\end{align*}
And for odd $n$, we have:
\begin{align*}
    J_{-n}(\beta)\sin((c-nm)t) + J_{n}(\beta)\sin((c+nm)t) \\
    + J_{-n}(\beta)\sin((-c-nm)t) + J_{n}(\beta)\sin((-c+nm)t) \\
     = -J_{n}(\beta)\sin((c-nm)t) + J_{n}(\beta)\sin((c+nm)t) \\
    - J_{n}(\beta)\sin(-1(c+nm)t) + J_{n}(\beta)\sin(-1(c-nm)t) \\
    = - J_{n}(\beta)\sin((c-nm)t) + J_{n}(\beta)\sin((c+nm)t) \\
    + J_{n}(\beta)\sin((c+nm)t) - J_{n}(\beta)\sin((c-nm)t) \\
    =\sin((c+nm)t)(J_{n}(\beta) + J_{n}(\beta)) + \sin((c-nm)t)(J_{-n}(\beta) + J_{-n}(\beta)) \\
    = 2J_{-n}(\beta)\sin((c+nm)t) + 2J_{n}(\beta)\sin((c+nm)t)
\end{align*}
As this must hold for all even and odd $n \in (-\infty, \infty)$, we can see that 
$$FM(A, c, m, \beta) + FM(A, -c, m, \beta) = 2A\sum_{n=-\infty, 2\nmid n}^\infty J_{n}(\beta)\sin((c + nm)t) $$
\end{proof}

\section*{An Algorithm to Produce an Arbitrary Frequency Using Non-Trivial Operators}
\label{sec:2_4}
 Despite the author's best efforts, they were unable to find an elegant way to do the algorithm with strictly simple FM operator pairs; however, by adding in amplitude modulated operators (AM), single-frequencies were found to have concise representations.\\

First, we define $AM(A, c, m) = \sin(m)\sin(c)$ to be a (phase-shifted) amplitude modulated signal with carrier frequency $c$ and modulating frequency $m$. Unlike how frequency modulated signals have infinitely many side bands, amplitude modulation always results in two side bands ($c-m, c+m$) which have the same magnitude as the carrier frequency. \\

In other words, 
$$\mathcal{F}(AM(A, c, m)) = A(\sin(c - m) + \sin(c) + \sin(c + m))$$

What this means is that we now have the ability to manipulate individual symmetric frequency bands from the FM operator through an AM operator. 

\begin{example} $f(t) = \sin(10\pi t)$
\end{example}

First, we start by defining the FM operator $FM_1(1,10\pi, m, \beta)$, where $m \in \mathbb{R}$ and $\beta : J_0(\beta) \neq \frac{1}{3} \in \mathbb{R}$.\\

Now, recall from Lemma 1, that if we add an FM operator with a negated modulating frequency, $FM_2(1, 10\pi, -m, \beta)$ we get a frequency spectrum consisting of only even harmonics from $FM_1$ with their amplitudes doubled. Recall also that these amplitudes are non-negative. \\

Hence, taking $FM_e = \frac{1}{2}(FM_1 + FM_2)$, we precisely get $FM_1$ without odd harmonics.\\

Now, notice that for each $n$, we have a symmetric side bands $$FM_{e,n} = J_{-n}(\beta)\sin((c - nm)t) + J_{n}(\beta)\sin((c+nm)t)$$

Now, define $AM_n(c, nm) = \sin(nm)\sin(c) = \sin((c - nm)t) + \sin(ct) + \sin((c+nm)t)$.\\

For each $n$, we can take $FM_e - J_n(\beta)AM_n(10\pi, m)$ to remove both side bands corresponding to the original index $n$ while subtracting $J_n(\beta)$ from the amplitude of the carrier. \\

Repeating for all $n \in [2, \infty) $, we are left with zero amplitudes for all $n \neq 0$, for which the carrier has amplitude $J_0(\beta) - J_2(\beta) - J_4(\beta) - ... - J_n(\beta) = J_0(\beta) - \sum_{n=1}^\infty J_{2n}(\beta)$.\\

Given $1 = \sum_{n=-\infty}^\infty J_n(\beta) = J_0 + \sum{n=1} 2J_{2n} \; \forall \beta$, we can write:
\begin{align*}
    1 &= J_0(\beta) + \sum_{n=1}^\infty 2J_{2n}(\beta) \\
    &= J_0(\beta) + 2\sum_{n=1}^\infty J_{2n}(\beta) \\
    \implies 2\sum_{n=1}^\infty J_{2n}(\beta) &= 1 - J_0(\beta)\\
    \implies \sum_{n=1}^\infty J_{2n}(\beta) &= \frac{1- J_0(\beta)}{2}
\end{align*}
Substituting back into the original expression, we have that the amplitude for $\sin(10\pi t)$ is:
\begin{align*}
    J_0(\beta) - \sum_{n=1}^\infty J_{2n}(\beta) &= J_0(\beta) - \frac{1- J_0(\beta)}{2}\\
    &= \frac{2J_0(\beta) - 1 + J_0(\beta)}{2}\\
    &= \frac{3J_0(\beta) - 1}{2}
\end{align*}

Therefore, given $\beta : 3J_0(\beta) \neq 1 \implies \beta \not\approx 1.81144$, if we multiply the result \\
$$FM_e - \sum_{n=1}^\infty J_{2n}(\beta)AM_{2n}(10\pi, nm)$$
by $\displaystyle\frac{2}{3J_0(\beta) - 1}$, we have:
\begin{align*}
    \frac{2}{3J_0(\beta) - 1}\left(FM_e - \sum_{n=1}^\infty J_{2n}(\beta)AM_{2n}(10\pi, nm)\right) &= \frac{2}{3J_0(\beta) - 1}\left(\frac{3J_0(\beta) - 1}{2}\sin(10\pi t)\right)\\
    &= \sin(10\pi t) 
\end{align*}

%%%%%%%%%%%%%%%%%%%%%%%%%%%%
% Bibliography
%%%%%%%%%%%%%%%%%%%%%%%%%%%%

% Force LaTeX to list the bibliography in the table of contents.

\clearpage
\addcontentsline{toc}{chapter}{\protect\numberline{}{Bibliography}}


% Now the actual bibliography:

\bibliography{bibfile}	
\bibliographystyle{alpha}



%%%%%%%%%%%%%%%%%%%%%%%%%%%%%%%%%%%%%%%%%%%
% Corrections
%
% This section should NOT be part of the
% original thesis.  It is only part of the
% corrected version, to be handed in by
% the second-to-last day of classes.
%
% So AFTER you hand in your thesis,
% un-comment the lines below and follow
% them with corrections, preferably in the
% style found in samplethesis.tex
%%%%%%%%%%%%%%%%%%%%%%%%%%%%%%%%%%%%%%%%%%%






\end{document}
